\documentclass[11pt,a4paper]{article}

\usepackage[utf8]{inputenc}
\usepackage[T1]{fontenc}
\usepackage{amsmath,amssymb,amsthm}
\usepackage{stmaryrd}
\usepackage{mathtools}
\usepackage{booktabs}
\usepackage{tabularx}
\usepackage{hyperref}
\usepackage[margin=2.5cm]{geometry}
\usepackage{natbib}
\usepackage{enumitem}
\usepackage{float}
\usepackage{xcolor}
\usepackage{tcolorbox}
\tcbuselibrary{skins,breakable}

\newtheorem{theorem}{Theorem}
\newtheorem{proposition}[theorem]{Proposition}
\newtheorem{corollary}[theorem]{Corollary}
\newtheorem{lemma}[theorem]{Lemma}
\newtheorem{definition}[theorem]{Definition}
\newtheorem{remark}[theorem]{Remark}
\newtheorem{example}[theorem]{Example}
\newtheorem{openquestion}{Open Question}
\newtheorem{axiom}{Axiom}

\newcolumntype{Y}{>{\raggedright\arraybackslash}X}

% -- Colors (from P9) --
\definecolor{rcolor}{HTML}{D85A30}
\definecolor{icolor}{HTML}{534AB7}
\definecolor{jcolor}{HTML}{1D9E75}
\definecolor{kcolor}{HTML}{185FA5}

\newcommand{\Op}{\mathcal{O}}
\newcommand{\rey}[1]{{\color{rcolor}#1}}
\newcommand{\imag}[1]{{\color{icolor}#1}}
\newcommand{\exac}[1]{{\color{jcolor}#1}}
\newcommand{\recu}[1]{{\color{kcolor}#1}}
\newcommand{\Real}{\operatorname{Re}}
\newcommand{\Imag}{\operatorname{Im}}
\newcommand{\SU}{\mathrm{SU}(2)}
\newcommand{\LZ}{\mathcal{Z}}

% -- tcolorbox styles --
\newtcolorbox{reinoR}{
  colback=rcolor!5, colframe=rcolor!40,
  fonttitle=\bfseries\color{rcolor!80!black},
  boxrule=0.4pt, arc=4pt,
  title={Axis $\rey{r}$ --- The Conceptual}
}
\newtcolorbox{reinoI}{
  colback=icolor!5, colframe=icolor!40,
  fonttitle=\bfseries\color{icolor!80!black},
  boxrule=0.4pt, arc=4pt,
  title={Axis $\imag{i}$ --- The Imaginary}
}
\newtcolorbox{reinoJ}{
  colback=jcolor!5, colframe=jcolor!40,
  fonttitle=\bfseries\color{jcolor!80!black},
  boxrule=0.4pt, arc=4pt,
  title={Axis $\exac{j}$ --- The Exact}
}
\newtcolorbox{reinoK}{
  colback=kcolor!5, colframe=kcolor!40,
  fonttitle=\bfseries\color{kcolor!80!black},
  boxrule=0.4pt, arc=4pt,
  title={Axis $\recu{k}$ --- The Recursive}
}
\newtcolorbox{pathologybox}[1][]{
  colback=black!3, colframe=black!20,
  fonttitle=\bfseries,
  boxrule=0.4pt, arc=4pt, #1
}
\newtcolorbox{axiombox}[1][]{
  colback=kcolor!3, colframe=kcolor!30,
  fonttitle=\bfseries\color{kcolor!80!black},
  boxrule=0.4pt, arc=4pt,
  title={#1}
}

\hypersetup{
  colorlinks=true,
  linkcolor=kcolor,
  citecolor=jcolor,
  urlcolor=rcolor,
  pdftitle={Quaternionic Logic: A G-Lattice Unifying Boolean and Fuzzy Frameworks},
  pdfauthor={J. Arturo Ornelas Brand},
  pdfsubject={Mathematical logic, many-valued logic, quaternionic algebra},
  pdfkeywords={quaternionic logic, G-lattice, residuated lattice, fuzzy logic, Boolean logic, SU(2), many-valued logic, recursive selection}
}
\pdfstringdefDisableCommands{%
  \def\rey#1{#1}%
  \def\imag#1{#1}%
  \def\exac#1{#1}%
  \def\recu#1{#1}%
}

\title{Quaternionic Logic: A G-Lattice Unifying\\Boolean and Fuzzy Frameworks}
\author{J.~Arturo Ornelas Brand\\
  \small Independent Researcher\\
  \small \texttt{arturoornelas62@gmail.com}}
\date{April 13, 2026}

\begin{document}
\maketitle

%======================================================================
\begin{abstract}
%======================================================================
A layered ontology of 72 semantic primitives across six algebraic
layers---Boolean, Fuzzy, Ordinal, Modal, Trivalent, and
Probabilistic---was previously shown to discriminate real knowledge
domains from pseudoscience controls. However, the six algebras
remained independent systems with no formal unification.

We construct a single quaternionic operator
$\Op=\{0,1,+,\imag{i},\exac{j},\recu{k}\}$ from which all six
algebras admit dimensional projections---four (Boolean, G\"odel
fuzzy, G\"odel--Kripke modal, and Kleene trivalent) as exact
recoveries, two (ordinal and probabilistic) as interpretive
correspondences. The operator
assigns to each concept a four-dimensional coordinate
$(\rey{r},\imag{i},\exac{j},\recu{k})\in[0,1]^4$, where $\rey{r}$
encodes crystallised truth, $\imag{i}$ encodes potentiality,
$\exac{j}$ encodes verifiability, and $\recu{k}$ encodes recursive
selection. Five axioms (K1--K5) characterise $\recu{k}$, and its
\emph{irreducibility} is proven: no function of
$(\rey{r},\imag{i},\exac{j})$ can reproduce $\recu{k}$.

Logical connectives on $V=[0,1]\times[-1,1]^3$ form a bounded
residuated lattice with involutive negation. An $\SU$ group acts on
truth values by conjugation, defining context transformations that
preserve truth while rotating epistemic character. The central result
is a \emph{commutation theorem}: lattice connectives commute with the
group action if and only if they depend exclusively on the real
component. Boolean logic emerges as the \emph{context-free}
(group-invariant) special case.

The combined structure is a \emph{G-lattice}: a residuated lattice
equipped with a non-abelian group action. Additional results include:
a cancellative--generative classification of nine algebras of
opposition, four diagnosable pathologies as axis collapses, and a
\emph{liquid zone} (all four axes positive) shown to be open, convex,
and path-connected. Completeness reduces to $G^\sim$ (G\"odel logic
with involutive negation). All results are verified computationally at
machine precision ($\sim10^{-16}$) over $50{,}000$~test cases.
\end{abstract}

%======================================================================
\section{Introduction}
\label{sec:intro}
%======================================================================

Six algebraic layers were introduced in \citet{p4} to organise the
72 semantic primitives of the prime-algebraic neurosymbolic
bridge~\citep{p1} into a dependency DAG: Boolean ($\{0,1\}$), Fuzzy
($\{0,+\}$), Ordinal ($\{<,=,>\}$), Modal ($\{\lozenge,\Box\}$),
Trivalent ($\{-,0,+\}$), and Probabilistic ($\{0,?\}$). Each layer
has its own formal logic, its own operators, and its own domain of
application. Yet they remained independent systems.

This paper asks two questions. First: \emph{is there a single
operator from which all six are recoverable as projections?}
Second: \emph{what formal logic does that operator support?}

We answer both affirmatively. The operator
$\Op=\{0,1,+,\imag{i},\exac{j},\recu{k}\}$ assigns to each
semantic state a point in a four-dimensional unit hypercube, where
each axis corresponds to an irreducible ontological category:
crystallised truth ($\rey{r}$), potentiality ($\imag{i}$),
verifiability ($\exac{j}$), and recursive selection ($\recu{k}$).
The six algebras are shadows cast by this four-dimensional object
onto lower-dimensional subspaces. The operator supports a residuated
lattice equipped with an $\SU$ group action---a structure we call a
\emph{G-lattice}---in which Boolean logic emerges as the unique
context-free special case.

\paragraph{Contributions.}
\begin{enumerate}[nosep]
\item Construction of a quaternionic operator from which all six
  algebraic layers are recoverable---four exactly (Boolean,
  G\"odel fuzzy, G\"odel--Kripke modal, Kleene trivalent) and
  two interpretively
  (Sections~\ref{sec:construction}--\ref{sec:recovery}).
\item Five axioms characterising the fourth axis $\recu{k}$, with
  a proof of its irreducibility
  (Sections~\ref{sec:axioms},~\ref{sec:irred}).
\item A bounded residuated lattice on $[0,1]\times[-1,1]^3$
  with exact recovery of four formal logics
  (Section~\ref{sec:lattice}).
\item A commutation theorem characterising Boolean logic as the
  group-invariant sublattice (Section~\ref{sec:glattice}).
\item Soundness and completeness, the latter reducing to
  $G^\sim$ (G\"odel logic with involutive negation) via a
  real-part independence lemma
  (Sections~\ref{sec:soundness},~\ref{sec:completeness}).
\item A cancellative--generative classification of nine algebras
  of opposition, with formal $\recu{k}$-embed\-dings for all five
  generative algebras
  (Section~\ref{sec:split}).
\item Four diagnosable pathologies as axis collapses, with a
  liquid zone shown to be open, convex, and path-connected
  (Section~\ref{sec:pathologies}).
\item A layer-based embedding of the prime algebra into the
  G-lattice domain, with a bilateral synthesis theorem
  characterising strict increase
  (Section~\ref{sec:embedding}).
\item A bilattice embedding of Belnap's $\mathbf{FOUR}$ into the
  G-lattice, preserving both truth and information orders as lattice
  homomorphisms (Section~\ref{sec:recovery}).
\item Computational verification of all results at machine precision
  (Section~\ref{sec:computation}).
\end{enumerate}

%======================================================================
\section{Preliminaries}
\label{sec:prelim}
%======================================================================

\subsection{Quaternions}

The quaternion algebra $\mathbb{H}$ is the four-dimensional real
algebra spanned by $\{1,i,j,k\}$ with multiplication rules
$i^2=j^2=k^2=ijk=-1$. Every quaternion $q=r+ai+bj+ck$ has real
part $\Real(q)=r$ and imaginary part $\Imag(q)=(a,b,c)$. The
conjugate is $\bar{q}=r-ai-bj-ck$, and for unit quaternions
($|q|=1$), $q^{-1}=\bar{q}$.

The group $\SU$ of unit quaternions acts on $\mathbb{H}$ by
conjugation: $q\cdot a = q\,a\,\bar{q}$. This action preserves
the real part, the imaginary norm, and the total norm:
\begin{align}
  \Real(q\,a\,\bar{q}) &= \Real(a), \label{eq:re-inv}\\
  \|\Imag(q\,a\,\bar{q})\| &= \|\Imag(a)\|, \label{eq:im-norm}\\
  \|q\,a\,\bar{q}\| &= \|a\|. \label{eq:total-norm}
\end{align}

\subsection{Residuated lattices}

A \emph{bounded residuated lattice} is a structure
$(L,\wedge,\vee,\otimes,\to,0,1)$ where $(L,\wedge,\vee,0,1)$ is
a bounded lattice, $\otimes$ is a monoidal operation (the t-norm),
and $\to$ is its residuum satisfying
$a\otimes b\leq c \iff a\leq b\to c$
\citep{galatos2007}.

%======================================================================
\section{Construction of the Operator}
\label{sec:construction}
%======================================================================

The construction proceeds in four steps, each adding one axis that no
previous combination could capture.

\subsection{Step 1: Recovering Boolean and Fuzzy --- the \texorpdfstring{$\rey{r}$}{r} axis}

The simplest attempt to subsume $\{0,1\}$ and $\{0,+\}$ is to take their
intersection: $\{0, +\}$, where $+$ abstracts over both $1$ (discrete)
and the open interval $(0,1)$ (continuous). But this loses information:
from $\{0,+\}$ one cannot recover which system was in use.

A true unifier must preserve both as special cases:
\begin{equation}
\Op_1 = \{0,\; \rey{1},\; \rey{+}\}
\end{equation}
where $0$ is absence, $\rey{1}$ is resolved (categorical) presence, and
$\rey{+}$ is partial (graded) presence. Recovery is immediate:
\begin{align}
\text{Boolean} &= \Op_1\big|_{\rey{+}=\varnothing} = \{0, \rey{1}\} \\
\text{Fuzzy}   &= \Op_1 \text{ with } \rey{+} \in (0,1),\; \rey{1} = \sup
\end{align}

This defines the first axis: $\rey{r} \in [0,1]$, the degree of
\emph{crystallised presence}. A concept with $\rey{r}=1$ is fully
resolved; with $\rey{r}=0.3$, partially so.

\subsection{Step 2: Incorporating modality --- the \texorpdfstring{$\imag{i}$}{i} axis}

The modal algebra $\{\lozenge, \Box\}$ distinguishes possibility from
necessity. Necessity ($\Box\,p$) maps to $\rey{r}=1$: the concept is
fully crystallised with no residual potential. But possibility
($\lozenge\,p$) introduces a genuinely new dimension: something that is
\emph{neither present nor absent} but \emph{latent}---potential that has
not yet actualised.

No value of $\rey{r}$ captures this. Adding a second axis:
\begin{equation}
\Op_2 = \{0,\; \rey{1},\; \rey{+},\; \imag{i}\}
\end{equation}
with the modal correspondence:
\begin{align}
\Box\,p   &\longleftrightarrow (\rey{r}=1,\; \imag{i}=0) \\
\lozenge\,p &\longleftrightarrow (\rey{r} \geq 0,\; \imag{i} > 0)
\end{align}

Both $\rey{r}$ and $\imag{i}$ admit continuous degrees, yielding a
\emph{semantic complex plane} where each state is a pair
$(\rey{r}, \imag{i}) \in [0,1]^2$.

\subsection{Step 3: The polarity axis --- \texorpdfstring{$\exac{j}$}{j}}

The trivalent algebra $\{-,0,+\}$ introduces \emph{polarity}: not just
presence or absence, but \emph{direction}. This is irreducible to
$\rey{r}$ (which has no sign) and to $\imag{i}$ (which has no
directionality).

\begin{equation}
\Op_3 = \{0,\; \rey{1},\; \rey{+},\; \imag{i},\; \exac{j}\}
\end{equation}

The $\exac{j}$ axis captures the degree of \emph{verifiability and
directional comparability}. It subsumes the ordinal algebra
$\{<,=,>\}$, which are relational operators that become possible only
when a measurement axis exists.

\subsection{Step 4: Recursive selection --- the \texorpdfstring{$\recu{k}$}{k} axis}

Three axes generate a static space: concepts have crystallisation,
potentiality, and directionality, but nothing \emph{moves}. The system
cannot prune, cannot advance, cannot select.

The probabilistic algebra $\{0,?\}$ of Layer~6 (Observer) hints at this
missing dimension. The ``$?$'' is not mere uncertainty---it is the state
of a system that \emph{looks at itself} and does not yet know what it
will find. This self-referential act---observation as recursive
selection---is irreducible to the other three axes.

The complete operator is:
\begin{equation}
\boxed{\;
\Op = \{0,\; \rey{1},\; \rey{+},\; \imag{i},\; \exac{j},\; \recu{k}\}
\;}
\end{equation}
where every semantic state lives in:
\begin{equation}
s = (\rey{r},\; \imag{i},\; \exac{j},\; \recu{k}) \in [0,1]^4
\end{equation}

The $\recu{k}$ axis measures the degree of \emph{recursive selection}:
how much a system prunes itself, rotates its perspective, and advances
irreversibly toward higher complexity. The four axes are irreducible:
no subset of three can generate the fourth. This mirrors the algebraic
structure of the quaternions $\mathbb{H}$, where $\{1,i,j,k\}$ are
the four basis elements and no three suffice to span the algebra. A
formal proof is given in Section~\ref{sec:irred}.

%======================================================================
\section{Axioms for Recursive Selection}
\label{sec:axioms}
%======================================================================

We propose five axioms that characterise $\recu{k}$.

\begin{axiombox}[Axiom K1 --- Range]
\begin{axiom}[Range]\label{ax:range}
$\recu{k} \in [0,1]$ for every semantic state. The value $k=0$
represents the complete absence of selection; $k=1$ represents maximal
selective activity.
\end{axiom}
\end{axiombox}

This is a structural axiom ensuring $\recu{k}$ lives in the same
bounded interval as the other axes, consistent with the hypercube
$H = [0,1]^4$.

\begin{axiombox}[Axiom K2 --- Stagnation Existence]
\begin{axiom}[Stagnation existence]\label{ax:stag}
The \emph{stagnation hyperplane}
\[
  \mathcal{S} = \{(r,i,j,0) \in H : r > 0,\; i > 0,\; j > 0\}
\]
is inhabited. That is, there exist semantic states with all three of
crystallisation, potentiality, and verifiability active, yet with zero
recursive selection.
\end{axiom}
\end{axiombox}

\noindent\textbf{Empirical grounding.} Consider Kodak in 1975: the
company had crystallised knowledge of photography ($\rey{r} > 0$),
imagined digital photography and built prototypes ($\imag{i} > 0$),
and verified feasibility with working systems ($\exac{j} > 0$)---yet
its membrane sealed and no selection occurred ($\recu{k} = 0$). The
company stagnated and eventually perished from conceptual heat death.

\begin{axiombox}[Axiom K3 --- Generation]
\begin{axiom}[Generation]\label{ax:gen}
The interaction of potentiality and verifiability \emph{enables} but
does not \emph{determine} selection:
\[
  \imag{i} > 0 \;\wedge\; \exac{j} > 0 \quad\Longrightarrow\quad
  \recu{k} \;\text{can be activated},
\]
but there exist states with $\imag{i} > 0$, $\exac{j} > 0$, and
$\recu{k} = 0$.
\end{axiom}
\end{axiombox}

This axiom formalises the quaternionic relation $\imag{i}\exac{j} =
\recu{k}$ as an \emph{enabling} relation rather than a deterministic
one. In standard quaternion algebra, $ij = k$ is an identity. Here, it
states that having both potentiality and verifiability is a
\emph{necessary precondition} for selection, but not sufficient---a
system may have both and still fail to select (Axiom~\ref{ax:stag}).

\begin{axiombox}[Axiom K4 --- Self-Inversion]
\begin{axiom}[Self-inversion]\label{ax:inv}
Recursive selection applied to itself produces inversion. In the
quaternion algebra $\mathbb{H}$, the basis element $k$ satisfies
$k^2=-1$: left-multiplication by $k$ twice reverses orientation.
\end{axiom}
\end{axiombox}

\noindent\textbf{Scope and domain.} K4 is a statement about the
algebraic structure of $\mathbb{H}$, not about the semantic
hypercube $H=[0,1]^4$: the map $-\mathrm{Id}:a\mapsto -a$ sends
$H$ outside itself (e.g., $-(0.5,0.3,0.2,0.1)\notin[0,1]^4$).
K4 therefore lives in $\mathbb{H}$ and constrains the
\emph{algebraic substrate} that the semantic axes mirror, not the
bounded domain of truth values.

\smallskip
\noindent\textbf{Left-multiplication vs.\ conjugation.}
The G-lattice (Section~\ref{sec:glattice}) uses conjugation
$q\cdot a=qa\bar{q}$, not left-multiplication $a\mapsto ka$.
These are distinct actions with different iteration behaviour:
\begin{itemize}[nosep]
\item \emph{Left-multiplication:} $(a\mapsto ka)\circ(a\mapsto ka)
  = (a\mapsto k^2 a) = (a\mapsto -a)$. Self-inversion.
\item \emph{Conjugation:} $(a\mapsto ka\bar{k})\circ
  (a\mapsto ka\bar{k}) = (a\mapsto k^2 a\bar{k}^2)
  = (a\mapsto(-1)a(-1)) = (a\mapsto a)$. Identity.
\end{itemize}
K4 captures the left-multiplication structure; the G-lattice
action captures the conjugation structure. This duality is not
an artefact of our construction: it is the standard distinction
in Clifford algebra $\mathrm{Cl}(0,2)\cong\mathbb{H}$ between
the left regular representation (spinor transformation) and the
twisted adjoint representation (vector
rotation)~\citep{lounesto2001,porteous1995}.
In the Clifford framework, $\mathbb{H}$ as an
$(\mathbb{H},\mathbb{H})$-bimodule encodes both actions: left
multiplication is the left module structure, and conjugation arises
from the bimodule structure via the canonical anti-involution.
The two actions coexist as structural features of the same algebra;
neither is derivable from the other.

\smallskip
\noindent\textbf{Semantic reading.} A system that selects
\emph{only which selections to make}---recursion without base
case---enters infinite regress and inverts. This is the algebraic
signature of Hofstadter's strange
loop~\citep{hofstadter1979}: self-reference that, taken to its
logical conclusion, flips the system.
Remark~\ref{rem:sign} below discusses how this quaternionic form
relates to the weaker involutive form found in categorical
realisations of $\recu{k}$.

\begin{axiombox}[Axiom K5 --- Activation Threshold]
\begin{axiom}[Activation threshold]\label{ax:thresh}
There exists a measurable function $\tau: (0,1]^3 \to [0,1]$ such
that the transition from $\recu{k} = 0$ to $\recu{k} > 0$ occurs when
the system's state $(\rey{r}, \imag{i}, \exac{j})$ crosses a critical
surface:
\[
  \recu{k} \;\text{activates when}\; \Psi(r, i, j) \geq \tau(r, i, j),
\]
where $\Psi$ is a system-dependent potential function.
\end{axiom}
\end{axiombox}

K5 is intentionally the weakest axiom: it postulates the
\emph{existence} of a threshold without specifying its form.
This generality is deliberate---different substrates have different
activation dynamics, and any stronger constraint (e.g., requiring
$\tau$ to be continuous, monotone, or given by a specific formula)
would restrict the framework to a particular substrate class.
The axiom is not vacuous: it constrains $\recu{k}$-activation to
depend on a \emph{deterministic} surface in $(r,i,j)$-space,
excluding stochastic or history-dependent activation. It is
non-trivially satisfiable: any continuous $\Psi$ with
$\Psi(r,i,j)=rij$ and any constant $\tau_0\in(0,1)$ produce a
concrete threshold surface $rij=\tau_0$.

\begin{remark}[Mutual consistency]
The five axioms are mutually consistent: K1 bounds $k$; K2 guarantees
stagnation exists; K3 says $i$ and $j$ enable but don't force $k$; K4
characterises pure $k$'s self-action; K5 says activation has a
threshold. No pair contradicts another, and K2's second clause is
implied by K3's second clause (making K2 the stronger existential
statement while K3 is the structural one).
\end{remark}

%======================================================================
\section{The Four Kingdoms}
\label{sec:kingdoms}
%======================================================================

Each axis corresponds to an ontological kingdom---an irreducible mode in
which information can exist.

\begin{reinoR}
What crystallises, fixes, defines. A concept with $\rey{r}=1$ is
categorical: triangle, prime number, conservation law. It has found its
form and holds it. The degree of resolved presence.
\emph{Ice.}
\end{reinoR}

\begin{reinoI}
What flows, invents, projects. A concept with $\imag{i}=1$ is pure
potentiality: the ``what if\ldots?'' that has not yet crystallised.
It is real, vivid, potent---but without form.
\emph{Plasma.}
\end{reinoI}

\begin{reinoJ}
What measures, compares, orders. A concept with $\exac{j}=1$ is fully
verifiable: the demonstration, the experiment, the measurement that
closes with \textit{Q.E.D.} and admits no debate. The scale without
temperature.
\emph{The thermometer.}
\end{reinoJ}

\begin{reinoK}
What selects, prunes, rotates, and advances. A concept with $\recu{k}=1$
is maximally recursive: it observes itself, discards what does not work,
and irreversibly increases complexity. It is the engine of evolution, the
phase transition, the moment where the system reorganises.
\emph{The spin.}
\end{reinoK}

\bigskip

\noindent
The four kingdoms are to information what the four fundamental forces are
to physics: irreducible, complementary, and jointly exhaustive. No
kingdom can be derived from the other three.

The state-of-matter metaphors---\emph{ice} for $\rey{r}$, \emph{plasma}
for $\imag{i}$, \emph{thermometer} for $\exac{j}$, \emph{spin} for
$\recu{k}$---are offered as \emph{mnemonic}, not definitional: they are
meant to make the ontological content easier to remember, not to ground
it. The definitional content is carried by the axioms of
Section~\ref{sec:axioms} and the construction of
Section~\ref{sec:construction}, not by the metaphors.

%======================================================================
\section{The Lattice of Connectives}
\label{sec:lattice}
%======================================================================

\begin{definition}[Truth values]
\label{def:truth}
A \emph{quaternionic truth value} is an element
$a=(\rey{r},\imag{i},\exac{j},\recu{k})\in V=[0,1]\times[-1,1]^3$.
The component $\rey{r}\in[0,1]$ is the \emph{degree of truth}.
The triple $(\imag{i},\exac{j},\recu{k})\in[-1,1]^3$ is the
\emph{epistemic character}.
\end{definition}

\begin{definition}[Connectives]
\label{def:connectives}
Define component-wise operations on $V=[0,1]\times[-1,1]^3$:
\begin{align}
  \mathrm{AND}(a,b) &= (\min(a_r,b_r),\;\min(a_i,b_i),\;\min(a_j,b_j),\;\min(a_k,b_k)), \\
  \mathrm{OR}(a,b) &= (\max(a_r,b_r),\;\max(a_i,b_i),\;\max(a_j,b_j),\;\max(a_k,b_k)), \\
  \mathrm{NOT}(a) &= (1-a_r,\;-a_i,\;-a_j,\;-a_k), \label{eq:not}\\
  (a\to b)_m &= \begin{cases}
    1 & \text{if } a_m\leq b_m,\\
    b_m & \text{otherwise,}
  \end{cases}
  \quad\text{for each component } m\in\{r,i,j,k\}. \label{eq:imp}
\end{align}
\end{definition}

\begin{proposition}[Lattice structure]
\label{prop:lattice}
$(V,\mathrm{AND},\mathrm{OR},\to,\mathbf{0},\mathbf{1})$
is a bounded residuated lattice with bounds
$\mathbf{0}=(0,-1,-1,-1)$ and $\mathbf{1}=(1,1,1,1)$
under the component-wise order.
\end{proposition}

\begin{proof}
Component-wise min and max on the bounded intervals
$[0,1]\times[-1,1]^3$ form a bounded distributive lattice with
$\mathbf{0}\leq a\leq\mathbf{1}$ for all $a\in V$.

\emph{Residuation.} We verify $\mathrm{AND}(a,b)\leq c$ iff
$a\leq(b\to c)$ for all $a,b,c\in V$, component by component.
Fix any component $m$.
$(\Rightarrow)$: If $b_m\leq c_m$, then $(b\to c)_m=1\geq a_m$.
If $b_m>c_m$, then $(b\to c)_m=c_m$ and we need $a_m\leq c_m$.
From $\min(a_m,b_m)\leq c_m$ and $b_m>c_m$ it follows that
$a_m\leq b_m$, hence $\min(a_m,b_m)=a_m\leq c_m$.
$(\Leftarrow)$: If $b_m\leq c_m$, then $\min(a_m,b_m)\leq b_m\leq c_m$.
If $b_m>c_m$, then $a_m\leq(b\to c)_m=c_m$, so
$\min(a_m,b_m)\leq a_m\leq c_m$.
\end{proof}

\begin{remark}[Involutive negation]
\label{rem:negation}
The negation~\eqref{eq:not} is \emph{not} the lattice negation
$a\to\mathbf{0}$ (which is G\"odel negation: $1$ if $a=\mathbf{0}$,
$\mathbf{0}$ otherwise). It is an additional \emph{involutive}
negation satisfying $\mathrm{NOT}(\mathrm{NOT}(a))=a$ and the
De Morgan laws:
$\mathrm{NOT}(\mathrm{AND}(a,b))=\mathrm{OR}(\mathrm{NOT}(a),\mathrm{NOT}(b))$,
since $1-\min(x,y)=\max(1-x,1-y)$ and $-\min(x,y)=\max(-x,-y)$.
The expanded structure is a residuated lattice with involutive
negation in the sense of \citet{esteva2000}.
\end{remark}

\begin{remark}[Product structure of the lattice]
\label{rem:product}
The lattice $(V,\mathrm{AND},\mathrm{OR},\to,\mathrm{NOT})$ is
isomorphic to the product of four independent chains:
$([0,1],\min,\max,\to_G,1{-}x)\times([-1,1],\min,\max,\to_G',{-}x)^3$.
In particular, the connectives do not couple the four components.
The quaternionic structure of the framework manifests not in the
lattice connectives but in two other places: (i)~the $\SU$ group
action (Section~\ref{sec:glattice}), which couples the imaginary
components via non-abelian rotation, and (ii)~the quaternionic
product on $S^3$ (Section~\ref{sec:quaternion}), which couples
all four components via Hamilton multiplication. A non-product
residuated lattice whose monoidal operation uses quaternion
multiplication directly is not available: quaternion multiplication
is not commutative, not idempotent, and does not preserve
$[0,1]^4$. The construction of a non-commutative residuated lattice
on $V$ using a quaternion-derived t-norm faces deep obstructions:
six families of lattice orders are incompatible with the Hamilton
product, and at most two of \{associativity, non-commutativity,
compatible lattice order\} can coexist on $V$---a theorem, not a
conjecture, proved via angular accumulation of iterated
products~\citep{p11_1}.
\end{remark}

%======================================================================
\section{Recovery of Formal Systems}
\label{sec:recovery}
%======================================================================

\begin{theorem}[Exact recovery]
\label{thm:recovery}
Under dimensional restrictions, the structure $(V,\mathrm{AND},\mathrm{OR},\mathrm{NOT},\to)$ reduces exactly to:
\begin{enumerate}[nosep]
\item \textbf{Boolean}: $i=j=k=0$, $r\in\{0,1\}$
  $\Rightarrow$ classical propositional logic.
\item \textbf{G\"odel fuzzy}: $i=j=k=0$, $r\in[0,1]$
  $\Rightarrow$ G\"odel fuzzy logic (min t-norm).
\end{enumerate}
\end{theorem}

\begin{proof}
Under restriction~(1), $V$ collapses to $\{0,1\}$: min is Boolean
AND, max is Boolean OR, $1-r$ is classical negation, and the
residuum gives material implication. Under~(2), $V$ collapses to
$[0,1]$ with the standard G\"odel connectives
\citep{hajek1998}. Both are isomorphisms, not merely embeddings.
\end{proof}

\subsection{Unified table of the six layers}

The six algebraic layers introduced in \citet{p4} are projections
of the operator $\Op$:

\begin{table}[H]
\centering
\renewcommand{\arraystretch}{1.4}
\caption{Six algebraic layers as projections of $\Op$.}
\label{tab:unified}
\begin{tabularx}{\textwidth}{@{}cl >{\raggedright}p{2cm} >{\raggedright}p{3cm} >{\raggedright\arraybackslash}p{4cm}@{}}
\toprule
\textbf{Layer} & \textbf{Name} & \textbf{Original algebra} &
\textbf{Unified operator} & \textbf{Active axes} \\
\midrule
1 & Point    & $\{0, 1\}$           & $\{0,\; \rey{1}\}$
  & $\rey{r}$ discrete \\
2 & Line     & $\{0, +\}$           & $\{0,\; \rey{+}\}$
  & $\rey{r}$ continuous \\
3 & Time     & $\{<, =, >\}$        & $\{\exac{<},\; \exac{=},\; \exac{>}\}$
  & $\exac{j}$ pure \\
4 & Plane    & $\{\lozenge, \Box\}$ & $\{\imag{i},\; \rey{1}\}$
  & $\rey{r} \times \imag{i}$ \\
5 & Volume   & $\{-, 0, +\}$        & $\{\exac{-j},\; 0,\; \exac{+j}\}$
  & $\exac{j^{\pm}}$ (polarity) \\
6 & Observer & $\{0, ?\}$           & $\{0,\;\rey{r}\imag{i}\exac{j}\recu{k}\}$
  & $\rey{r} \times \imag{i} \times \exac{j} \times \recu{k}$ \\
\bottomrule
\end{tabularx}
\end{table}

\noindent
The layers progressively activate axes: the first two operate on
$\rey{r}$ alone (presence), the middle layers activate one or two
additional axes ($\imag{i}$, $\exac{j}$), and Layer~6 (Observer) is
the only one inhabiting all four dimensions simultaneously. Layers
1--5 set $\recu{k}=0$ (they predate selection).

\begin{theorem}[Modal and trivalent recovery]
\label{thm:modal_trivalent}
Under dimensional restrictions, two further logics are exactly
recoverable:
\begin{enumerate}[nosep,start=3]
\item \textbf{Modal (G\"odel--Kripke)}: Restrict to $j=k=0$,
  $r,i\in[0,1]$, and equip the G\"odel restriction with Kripke
  frames valued in $[0,1]$, reading $\Box\varphi$ as the infimum
  of $\varphi$ over accessible worlds and $\lozenge\varphi$ as
  the supremum. Then the $\{0,1\}$-subalgebra recovers classical
  modal logic exactly. Completeness of G\"odel modal logics
  (including S4 and S5 at the Boolean restriction) is established
  by \citet{caicedo2010}; the general construction for residuated
  lattices is due to \citet{bou2011} and \citet{fitting1991}.
\item \textbf{Trivalent (Kleene)}: Restrict to $i=k=0$,
  $j\in\{-1,0,+1\}$. The connectives $\min$, $\max$, and
  $-j$ on the three-element chain $\{-1,0,+1\}$ are exactly
  the strong Kleene truth tables (with the identification
  $-1=\mathrm{false}$, $0=\mathrm{unknown}$,
  $+1=\mathrm{true}$). The restriction is closed under all three
  operations and constitutes an isomorphism, not an
  approximation. More generally, {\L}ukasiewicz $n$-valued logic
  is exactly the $n$-element subalgebra of the standard
  MV-algebra on $[0,1]$~\citep{cignoli2000}.
\end{enumerate}
\end{theorem}

\begin{proof}
\emph{(3) Modal.} The G\"odel chain $([0,1],\min,\max,\to_G)$ is
a Heyting algebra and hence admits Kripke semantics. By
\citet{caicedo2010}, Theorem~3.8, the standard G\"odel modal logic
over $[0,1]$ is complete with respect to Kripke frames with
$[0,1]$-valued accessibility. When the valuation is restricted to
$\{0,1\}$, the accessibility relation becomes crisp, the G\"odel
residuum becomes material implication, and the resulting logic is
classical modal logic (K, S4, or S5 depending on frame conditions).

\emph{(4) Trivalent.} On $\{-1,0,+1\}$ with the natural order
$-1<0<+1$: $\min$ is Kleene conjunction, $\max$ is Kleene
disjunction, and $x\mapsto -x$ is Kleene negation. These are
exactly the strong Kleene truth tables. The subset is closed under
all three operations (trivially: $\min$, $\max$, and negation of
elements of $\{-1,0,+1\}$ remain in $\{-1,0,+1\}$). This gives
an isomorphism between the restricted sublattice and
$\mathrm{K}_3$.
\end{proof}

\begin{remark}[Remaining interpretive correspondences]
\label{rem:interpretive}
Two layers remain as interpretive correspondences rather than
exact isomorphisms:
\begin{enumerate}[nosep,start=5]
\item \textbf{Ordinal}: $j=k=0$, $i$ encoding rank.
\item \textbf{Probabilistic}: $j=k=0$, $i$ encoding variance.
\end{enumerate}
These are suggestive readings of the restricted subspaces. Ordinal
and probabilistic logics have richer structure (well-orders,
$\sigma$-algebras) than the quaternionic restriction provides.
\end{remark}

\subsection{Belnap--Dunn bilattice embedding}

Belnap's four-valued logic~\citep{belnap1977} uses the bilattice
$\mathbf{FOUR} = \{\bot, \mathbf{f}, \mathbf{t}, \top\}$ with two
independent lattice orders: a \emph{truth} order
($\mathbf{f} \leq_t \bot \leq_t \mathbf{t}$,
$\mathbf{f} \leq_t \top \leq_t \mathbf{t}$, with $\bot$ and $\top$
incomparable) and an \emph{information} order
($\bot \leq_i \mathbf{f} \leq_i \top$,
$\bot \leq_i \mathbf{t} \leq_i \top$, with $\mathbf{f}$ and
$\mathbf{t}$ incomparable).

\begin{theorem}[Bilattice embedding]
\label{thm:belnap}
$\mathbf{FOUR}$ embeds into $V = [0,1] \times [-1,1]^3$ as a
simultaneous lattice homomorphism for both orders. Specifically,
there exists $\varphi: \mathbf{FOUR} \to V$ such that:
\begin{enumerate}[nosep]
  \item $\varphi$ preserves the truth order as the component-wise
    order on $V$: $a \leq_t b$ iff $\varphi(a) \leq_{\mathrm{cw}}
    \varphi(b)$.
  \item $\varphi$ preserves the information order as the
    component-wise order on absolute imaginary magnitudes:
    $a \leq_i b$ iff $|\mathrm{Im}(\varphi(a))| \leq_{\mathrm{cw}}
    |\mathrm{Im}(\varphi(b))|$.
  \item Both truth and information lattice operations (meet and join)
    are preserved.
\end{enumerate}
\end{theorem}

\begin{proof}
Choose parameters $r_0 < r_1$ in $[0,1]$, $\, s < 0 < t$ in
$[-1,1]$ with $|s| > t$, and $0 < d$ in $[0,1]$. Define:
\begin{align*}
\varphi(\bot) &= (r_0,\; t,\; 0,\; 0), &
\varphi(\mathbf{f}) &= (r_0,\; s,\; 0,\; 0), \\
\varphi(\mathbf{t}) &= (r_1,\; t,\; d,\; 0), &
\varphi(\top) &= (r_1,\; s,\; d,\; 0).
\end{align*}
The construction assigns each coordinate to either case~(A)
($\bot_x = \mathbf{f}_x$, $\top_x = \mathbf{t}_x$) or case~(B)
($\bot_x = \mathbf{t}_x$, $\top_x = \mathbf{f}_x$):
\emph{$r$}~takes case~(A), \emph{$i$}~takes case~(B), and
\emph{$j$}~takes case~(A).

\smallskip
\emph{(1) Truth order.} Since $s < t$ and $r_0 < r_1$, the
component-wise relations are:
$\varphi(\mathbf{f}) \leq \varphi(\bot)$ (only $i$ changes:
$s \leq t$),
$\varphi(\bot) \leq \varphi(\mathbf{t})$ ($r_0 \leq r_1$,
$0 \leq d$), and similarly for $\varphi(\mathbf{f}) \leq
\varphi(\top) \leq \varphi(\mathbf{t})$. For incomparability of
$\bot$ and $\top$: coordinate $i$ gives $t > s$ ($\bot$ higher)
while $r$ gives $r_0 < r_1$ ($\top$ higher).

\smallskip
\emph{(2) Information order.} The absolute imaginary magnitudes
are $|\mathrm{Im}(\varphi(\bot))| = (t, 0)$,
$|\mathrm{Im}(\varphi(\mathbf{f}))| = (|s|, 0)$,
$|\mathrm{Im}(\varphi(\mathbf{t}))| = (t, d)$,
$|\mathrm{Im}(\varphi(\top))| = (|s|, d)$.
Since $|s| > t > 0$ and $d > 0$: $\bot$ has the smallest
magnitudes, $\top$ the largest, and $\mathbf{f}$ (large $|i|$,
no $|j|$) is incomparable with $\mathbf{t}$ (small $|i|$,
large $|j|$).

\smallskip
\emph{(3) Lattice operations.} By construction,
$\max_{\mathrm{cw}}(\varphi(\bot), \varphi(\top)) =
(r_1, t, d, 0) = \varphi(\mathbf{t})$ and
$\min_{\mathrm{cw}}(\varphi(\bot), \varphi(\top)) =
(r_0, s, 0, 0) = \varphi(\mathbf{f})$, matching the truth lattice.
Similarly,
$\max_{|\mathrm{Im}|}(\varphi(\mathbf{f}), \varphi(\mathbf{t})) =
(|s|, d) = |\mathrm{Im}(\varphi(\top))|$ and
$\min_{|\mathrm{Im}|}(\varphi(\mathbf{f}), \varphi(\mathbf{t})) =
(t, 0) = |\mathrm{Im}(\varphi(\bot))|$, matching the information
lattice.
\end{proof}

\begin{remark}[Sign-magnitude decoupling]
\label{rem:belnap-decoupling}
The embedding exploits the sign degree of freedom in the imaginary
components: a \emph{negative} imaginary coordinate is \emph{below}
its positive counterpart in component-wise (truth) order, but its
\emph{absolute value} is \emph{above}---decoupling truth from
information. This is the algebraic mechanism by which the G-lattice
accommodates two independent orders on the same space. In Belnap's
semantics, $\mathbf{f}$ and $\mathbf{t}$ carry equal information
($|\mathrm{Im}|$ concentrated on different axes) but opposite truth;
$\bot$ carries minimal information while $\top$ carries maximal
information, both at neutral truth.
\end{remark}

%======================================================================
\section{The G-Lattice: Context as Group Action}
\label{sec:glattice}
%======================================================================

\begin{definition}[Context transformation]
\label{def:context}
A \emph{context} is a unit quaternion $q\in\SU$. The context
transformation of a truth value $a$ is the conjugation action:
\[
  q\cdot a \;=\; q\,a\,\bar{q}.
\]
\end{definition}

\begin{definition}[G-lattice]
\label{def:glattice}
The \emph{quaternionic G-lattice} is the triple $(L,G,\cdot)$ where:
\begin{itemize}[nosep]
\item $L=(V,\mathrm{AND},\mathrm{OR},\to,\mathbf{0},\mathbf{1})$ is the
  residuated lattice from Proposition~\ref{prop:lattice}, equipped
  with the involutive negation $\mathrm{NOT}$ from
  Remark~\ref{rem:negation},
\item $G=\SU$ is the group of unit quaternions,
\item $\cdot$ is the conjugation action from Definition~\ref{def:context}.
\end{itemize}
\end{definition}

\begin{proposition}[Action properties]
\label{prop:action}
The conjugation action satisfies:
\begin{enumerate}[nosep]
\item $\Real(q\cdot a)=\Real(a)$ \hfill(truth is objective),
\item $\|\Imag(q\cdot a)\|=\|\Imag(a)\|$ \hfill(epistemic weight is conserved),
\item $\|q\cdot a\|=\|a\|$ \hfill(information is conserved),
\item $q\cdot(a\cdot_{\mathbb{H}} b)=(q\cdot a)\cdot_{\mathbb{H}}(q\cdot b)$ \hfill(algebra automorphism).
\end{enumerate}
\end{proposition}

\begin{proof}
Properties (1)--(3) follow from equations~\eqref{eq:re-inv}--\eqref{eq:total-norm}.
Property (4): $q(ab)\bar{q}=(qa\bar{q})(qb\bar{q})$ by
associativity and $\bar{q}q=1$.
\end{proof}

\begin{proposition}[Non-closure of the action on $V$]
\label{prop:nonclosure}
The conjugation action is not closed on $V=[0,1]\times[-1,1]^3$.
That is, there exist $a\in V$ and $q\in\SU$ such that
$q\cdot a\notin V$.
\end{proposition}

\begin{proof}
Let $a=(0.5,\;1,\;1,\;0)\in V$. Then
$\|\Imag(a)\|=\sqrt{1^2+1^2+0^2}=\sqrt{2}$.
Conjugation acts on $\Imag(a)$ as an $\mathrm{SO}(3)$ rotation
preserving $\|\Imag\|=\sqrt{2}$. Choose $q\in\SU$ corresponding
to the $\mathrm{SO}(3)$ rotation that maps $(1,1,0)$ to
$(0,0,\sqrt{2})$. Then $q\cdot a=(0.5,\;0,\;0,\;\sqrt{2})$.
Since $\sqrt{2}>1$, the $\recu{k}$-component exceeds $[-1,1]$,
so $q\cdot a\notin V$.
\end{proof}

\begin{proposition}[No domain is simultaneously closed under
connectives and the group action]
\label{prop:no_joint_domain}
There is no subset $D\subseteq[0,1]\times\mathbb{R}^3$ that is
simultaneously closed under $\mathrm{OR}$ (component-wise max)
and under $\SU$ conjugation, unless $D$ is contained in
$\{a:\Imag(a)=0\}$.
\end{proposition}

\begin{proof}
Suppose $D$ is closed under both operations and contains some
$a$ with $\Imag(a)\neq0$. By $\SU$-closure, $D$ contains all
rotations of $\Imag(a)$; in particular, it contains elements
$a_1,a_2$ with $\Real(a_1)=\Real(a_2)$ and disjoint support in
the imaginary components (e.g., $a_1=(r,\alpha,0,0)$ and
$a_2=(r,0,\alpha,0)$ for $\alpha=\|\Imag(a)\|>0$). Then
$\mathrm{OR}(a_1,a_2)=(r,\alpha,\alpha,0)$ with
$\|\Imag\|=\alpha\sqrt{2}>\alpha$. Iterating: $\mathrm{OR}$
with a third rotation gives $\|\Imag\|=\alpha\sqrt{3}$, and
further rotations and ORs produce arbitrarily large
$\|\Imag\|$. Hence $D$ is unbounded in the imaginary
components, contradicting any finite bound. The only escape is
$\alpha=0$, i.e., $\Imag(a)=0$.
\end{proof}

\begin{remark}[Domain separation]
\label{rem:domain}
Propositions~\ref{prop:nonclosure} and~\ref{prop:no_joint_domain}
establish that the lattice and the group action necessarily live on
different domains: the lattice on $V=[0,1]\times[-1,1]^3$, the
group action on the ambient space
$W=[0,1]\times\mathbb{R}^3\supseteq V$. This is not a deficiency
but a structural fact: the lattice connectives (min, max, G\"odel
residuum, involutive negation) extend naturally to $W$, and the
real-part independence lemma (Lemma~\ref{lem:re-indep}) holds
identically on $W$. Therefore all logical results---tautologies,
soundness, completeness, and the commutation theorem---are
unaffected by the domain extension. The two structures interact
through $\Real$, which both preserve.
\end{remark}

\begin{theorem}[Commutation theorem]
\label{thm:commutation}
Let $f:V^n\to V$ be any connective (AND, OR, NOT, $\to$, or a
composition thereof). Then
\[
  q\cdot f(a_1,\ldots,a_n)=f(q\cdot a_1,\ldots,q\cdot a_n)
  \quad\text{for all } q\in\SU
\]
if and only if $f$ depends only on the real components
$(\Real(a_1),\ldots,\Real(a_n))$.
\end{theorem}

\begin{proof}
$(\Leftarrow)$ If $f$ depends only on real parts, then
$\Real(q\cdot a_m)=\Real(a_m)$ by Proposition~\ref{prop:action}(1),
so $f$ receives the same inputs regardless of $q$. The output has
zero imaginary part (since $f$ outputs only a real-part result),
and conjugation fixes real quaternions. Hence both sides are equal.

$(\Rightarrow)$ Suppose $f$ depends on some imaginary component.
Choose $a$ such that the imaginary part of $f(a)$ is nonzero. Then
there exists $q\in\SU$ that rotates $\Imag(f(a))$ to a different
direction, while $f(q\cdot a)$ computes min/max on the rotated
components of $a$---which generically differs from the rotation of
min/max of the original components (since min does not commute with
rotation). Hence the two sides differ.
\end{proof}

\begin{example}[Explicit non-commutation]
\label{ex:noncommute}
Let $a=(0.7,0.3,0.5,0.1)$, $b=(0.4,0.8,0.2,0.6)$, and
$q=\tfrac{1}{\sqrt{2}}(1,1,0,0)$.
Then $q\cdot\mathrm{AND}(a,b)=(0.4,\;0.3,\;{-0.1},\;0.2)$ while
$\mathrm{AND}(q\cdot a,\;q\cdot b)=(0.4,\;0.3,\;{-0.6},\;0.2)$;
the $j$-components differ ($-0.1$ vs $-0.6$).
The real parts agree ($0.4=0.4$), confirming that the
non-commutation is purely epistemic.
\end{example}

\begin{corollary}[Boolean logic is context-free]
\label{cor:boolean}
The sublattice $\{a\in V:\Imag(a)=0\}$ is the unique maximal
sublattice on which all connectives commute with every context
transformation. This is the Boolean/fuzzy restriction.
\end{corollary}

\begin{definition}[Context sensitivity]
\label{def:sensitivity}
The \emph{context sensitivity} of a truth value $a$ is
\[
  \sigma(a) = 2\,\|\Imag(a)\|.
\]
This equals the maximum displacement $\max_{q\in\SU}\|q\cdot a - a\|$.
\end{definition}

\begin{remark}
The context sensitivity provides a continuous spectrum:
$\sigma=0$ corresponds to Boolean logic (context-free),
and increasing $\sigma$ gives progressively more
context-dependent logics. For truth values in $V$, the spectrum
is parameterised by $\|\Imag(a)\|\in[0,\sqrt{3}]$.
\end{remark}

%======================================================================
\section{Soundness}
\label{sec:soundness}
%======================================================================

\begin{definition}[Truth predicate]
\label{def:true}
For a threshold $\tau\in[0,1]$, a truth value $a$ is \emph{true}
if $\Real(a)\geq\tau$.
\end{definition}

The inference rules of the quaternionic G-lattice are:

\begin{enumerate}[nosep,label=\textbf{R\arabic*}.]
\item \textbf{Modus ponens:} From $a$ and $a\to b$, infer $b$.
\item \textbf{Context shift:} From $a$, infer $q\cdot a$ for any $q\in\SU$.
\item \textbf{Context composition:} $q_2\cdot(q_1\cdot a)=(q_2 q_1)\cdot a$.
\end{enumerate}

\begin{theorem}[Soundness]
\label{thm:soundness}
The inference rules \textbf{R1}--\textbf{R3} are sound with respect
to the truth predicate of Definition~\ref{def:true}.
\end{theorem}

\begin{proof}
\textbf{R1} (Modus ponens): If $\Real(a)\geq\tau$ and
$\Real(a\to b)\geq\tau$, then either $a_r\leq b_r$ (in which case
$b_r\geq a_r\geq\tau$) or $a_r>b_r$ (in which case
$\Real(a\to b)=b_r\geq\tau$). In both cases $\Real(b)\geq\tau$.
This is the standard soundness argument for the G\"odel residuum
\citep{hajek1998,galatos2007}.

\textbf{R2} (Context shift): By Proposition~\ref{prop:action}(1),
$\Real(q\cdot a)=\Real(a)\geq\tau$.

\textbf{R3} (Context composition): By associativity of quaternion
multiplication: $q_2(q_1 a\bar{q}_1)\bar{q}_2=(q_2 q_1)a(\overline{q_2 q_1})$.
\end{proof}

\begin{remark}
Soundness decomposes cleanly: \textbf{R1} is inherited from the
residuated lattice literature; \textbf{R2} follows from the
real-part invariance of conjugation; \textbf{R3} is the group law.
No new proof technique is required---the G-lattice soundness is the
\emph{product} of two well-understood soundness results.
\end{remark}

%======================================================================
\section{Completeness}
\label{sec:completeness}
%======================================================================

\begin{lemma}[Real-part independence]
\label{lem:re-indep}
For every propositional formula $\varphi$ built from AND, OR, NOT,
$\to$ and every valuation $v:\mathrm{Atoms}\to V$,
\[
  \Real(\llbracket\varphi\rrbracket_v)
  \;\text{depends only on}\;
  (\Real(v(p_1)),\ldots,\Real(v(p_n))).
\]
That is, the imaginary components of the atoms do not affect the
real part of the result.
\end{lemma}

\begin{proof}
By structural induction on $\varphi$. For each connective:
\begin{align*}
  \Real(\mathrm{AND}(a,b)) &= \min(\Real(a),\Real(b)), \\
  \Real(\mathrm{OR}(a,b))  &= \max(\Real(a),\Real(b)), \\
  \Real(\mathrm{NOT}(a))   &= 1 - \Real(a), \\
  \Real(a\to b)             &= \begin{cases}1 & \text{if }\Real(a)\leq\Real(b),\\ \Real(b) & \text{otherwise.}\end{cases}
\end{align*}
In each case the output's real part is a function of the inputs'
real parts alone. The induction step follows by compositionality.
\end{proof}

\begin{theorem}[Completeness]
\label{thm:completeness}
Let $\varphi$ be a propositional formula built from
$\mathrm{AND},\mathrm{OR},\mathrm{NOT},\to$.
\begin{enumerate}[nosep]
\item $\varphi$ is a G-lattice tautology
  ($\Real(\llbracket\varphi\rrbracket_v)\geq\tau$ for every
  $v:\mathrm{Atoms}\to V$)
  if and only if $\varphi$ is a tautology of the $[0,1]$-algebra
  $(\min,\max,1{-}x,\to_G)$, where $\to_G$ is the G\"odel
  residuum.
\item For the negation-free fragment
  ($\varphi$ built from $\mathrm{AND},\mathrm{OR},\to$ only),
  this reduces to G\"odel--Dummett completeness
  \citep{dummett1959,hajek1998}.
\item For the full fragment (with $\mathrm{NOT}$), the logic is
  $G^\sim$---G\"odel logic expanded with an involutive negation---which
  is complete by \citet{esteva2000}.
\end{enumerate}
\end{theorem}

\begin{proof}
By Lemma~\ref{lem:re-indep}, $\Real(\llbracket\varphi\rrbracket_v)$
depends only on the real parts of the atomic valuations. Therefore
$\varphi$ is a G-lattice tautology iff it is a tautology of the
real-part projection, which is the $[0,1]$-algebra with $\min$,
$\max$, the G\"odel residuum, and the involutive negation $1-x$.

For the negation-free fragment, this algebra is the standard G\"odel
chain, and completeness follows from \citet{dummett1959,hajek1998}.
For the full fragment, the algebra is a G\"odel algebra expanded with
involutive negation; completeness of this class ($G^\sim$) is
established in \citet{esteva2000}.
\end{proof}

\begin{remark}
The involutive negation $\mathrm{NOT}(a)=1-a$ is \emph{not} the
G\"odel negation $a\to\mathbf{0}$. In particular, $\mathrm{NOT}(\mathrm{NOT}(a))=a$ is a
tautology (double negation holds), while
$(\varphi\to\psi)\to(\mathrm{NOT}(\psi)\to\mathrm{NOT}(\varphi))$
(contraposition) is not. These are characteristic of $G^\sim$,
not of standard G\"odel logic.

The imaginary dimensions enrich the logic with context sensitivity,
epistemic character, and path-dependent reasoning, but they do not
alter which formulas are tautologies. The ``new'' content
of the quaternionic G-lattice is not in its tautologies but in the
\emph{structure around} them: the group action, the context spectrum,
and the cause--effect chains.
\end{remark}

%======================================================================
\section{Irreducibility of \texorpdfstring{$\recu{k}$}{k}}
\label{sec:irred}
%======================================================================

We first establish a key structural property.

\begin{proposition}[Non-injectivity of projection]
\label{prop:noninj}
The projection $\pi: H \to [0,1]^3$ defined by $\pi(r,i,j,k) =
(r,i,j)$ is not injective on the inhabited region of~$H$. That is,
there exist states $s_0, s_1 \in H$ with $\pi(s_0) = \pi(s_1)$ and
$s_0 \neq s_1$.
\end{proposition}

\begin{proof}
By Axiom~\ref{ax:stag} (K2), there exists a stagnation state $s_0 =
(r_0, i_0, j_0, 0)$ with $r_0, i_0, j_0 > 0$. The liquid-zone state
$s_1 = (r_0, i_0, j_0, k_1)$ with any $k_1 > 0$ also lies in $H$
(since $[0,1]^4$ contains all such points). Both states project to
$(r_0, i_0, j_0)$ under $\pi$, but $s_0 \neq s_1$ since their
$\recu{k}$-components differ.
\end{proof}

\begin{theorem}[Irreducibility of $\recu{k}$]
\label{thm:irred}
There exists no function $f: [0,1]^3 \to [0,1]$ such that
$k = f(r, i, j)$ for all semantic states $(r, i, j, k) \in H$.
\end{theorem}

\begin{proof}
Suppose, for contradiction, that such a function $f$ exists. Let
$s_0 = (r_0, i_0, j_0, 0) \in \mathcal{S}$ be a stagnation state
(exists by Axiom~\ref{ax:stag}). Then $f(r_0, i_0, j_0) = 0$.

Now consider the state $s_1 = (r_0, i_0, j_0, k_1)$ with $k_1 > 0$
(which exists in $H$ by Proposition~\ref{prop:noninj}). Since $f$
determines $k$ from $(r,i,j)$, we must have $f(r_0, i_0, j_0) = k_1 >
0$.

But $f(r_0, i_0, j_0) = 0$ (from $s_0$) and $f(r_0, i_0, j_0) = k_1
> 0$ (from $s_1$) is a contradiction, since $f$ is a function and must
assign a unique value to each input. Therefore no such $f$ exists.
\end{proof}

\begin{corollary}
\label{cor:independence}
The four axes $(\rey{r}, \imag{i}, \exac{j}, \recu{k})$ are
algebraically independent in the sense that no axis can be expressed as
a function of the remaining three.
\end{corollary}

\begin{proof}
The irreducibility of $\recu{k}$ with respect to $(\rey{r}, \imag{i},
\exac{j})$ is Theorem~\ref{thm:irred}. The irreducibility of each of
$\rey{r}$, $\imag{i}$, $\exac{j}$ with respect to the other three
follows by an identical argument using the corresponding pathologies:
dogma ($\imag{i} = 0$ with $r, j, k > 0$), delirium ($\rey{r} = 0$
with $i, j, k > 0$), and pseudoscience ($\exac{j} = 0$ with $r, i, k
> 0$). Each pathology provides a state where one axis is zero while
the others are positive, and the liquid zone provides states where all
four are positive with the same values on the three non-zero axes.
\end{proof}

\begin{remark}[Strength of the proof]
The proof requires only Axiom K2 (stagnation existence) and the
geometric fact that the hypercube $[0,1]^4$ contains both stagnation
and liquid-zone states sharing the same $(r,i,j)$ projection. No
dynamical assumptions, no quaternionic structure, and no threshold
mechanism are needed.
\end{remark}

%======================================================================
\section{The Cancellative--Generative Classification}
\label{sec:split}
%======================================================================

Nine algebraic structures of opposition arise naturally in the
prime-algebraic framework. We classify them according to whether they
require $\recu{k}$.

\begin{definition}[Cancellative and generative algebras]
An algebra of opposition is \emph{cancellative} if combining an
element with its opposite produces a neutral element (zero, one, or
identity matrix). It is \emph{generative} if the combination produces
structure that did not exist in either operand.
\end{definition}

\begin{theorem}[Cancellative--generative classification]
\label{thm:split}
The nine algebras of opposition split into two classes:

\smallskip
\noindent\textbf{Cancellative} (operate within $(\rey{r}, \imag{i}, \exac{j})$, do not require $\recu{k}$):
\begin{enumerate}[nosep]
  \item Additive inverse: $a + (-a) = 0$
  \item Multiplicative inverse: $a \cdot (1/a) = 1$
  \item Cyclic group $\mathbb{Z}/n\mathbb{Z}$: $a + (n-a) \equiv 0$
  \item Rotation $\mathrm{SO}(2)$: $R(\theta) \cdot R(-\theta) = I$
\end{enumerate}

\smallskip
\noindent\textbf{Generative} (require $\recu{k}$):
\begin{enumerate}[nosep, start=5]
  \item Categorical duality: $(-)^{\mathrm{op}}$ reverses all arrows
  \item Kernel/Image: $\dim(\ker f) + \dim(\mathrm{im}\, f) = \dim V$
  \item Eigenvalue spectrum: $Av = \lambda v$ classifies modes
  \item Brouwer fixed point: $\exists\, x^*: f(x^*) = x^*$
  \item Homology/Betti numbers: structure of absence
\end{enumerate}
\end{theorem}

\begin{proof}
\textbf{Cancellative algebras (1--4).} Each has the form $a \ast
a^{-1} = e$ where $e$ is an identity element and $\ast$ is the group
operation. This is formally verifiable: given $a$, the result $e$ is
uniquely determined by the group axioms---no selection is needed.
These algebras operate by producing a neutral element, which is a
fixed point of the operation. They model opposition that
\emph{annihilates}: wave plus anti-wave yields silence, credit plus
debt yields zero. All four are expressible entirely in terms of the
content axes $(\rey{r}, \imag{i}, \exac{j})$: additive inverse
operates on $\rey{r}$ (cancellation of crystallised content),
multiplicative inverse on $\exac{j}$ (normalisation of measurement),
cyclic return on $\imag{i}$ (return of potentiality), and rotation on
the $(\rey{r}, \imag{i})$ plane (phase shift without selection).

\textbf{Generative algebras (5--9).}
The classification criterion for the generative class is the
shared structural property: each involves a
\emph{meta-operation}---an operation that acts on the structure
of a system rather than within it.
Theorems~\ref{thm:catdual}--\ref{thm:homology} below prove
individually that each algebra's characteristic operation satisfies
the three defining properties of $\recu{k}$: (a)~meta-operation
on structure, (b)~generative, and (c)~self-inverting.

\emph{(5) Categorical duality} reverses \emph{all} arrows of a
diagram simultaneously---a global structural transformation, not a
local computation. It produces ``free theorems'': any result about
products has a dual about coproducts. This requires operating
\emph{on} the structure of $(r,i,j)$, not \emph{within} it.

\emph{(6) Kernel/Image} detects what a transformation
\emph{destroyed} versus what it \emph{created}. This is a
meta-diagnostic: the system auditing its own information loss.

\emph{(7) Eigenvalue spectrum} classifies the \emph{modes} of action
($\lambda > 0$: expansion, $\lambda < 0$: reflection, $\lambda \in
\mathbb{C}$: oscillation, $\lambda = 0$: annihilation). Classifying
modes requires observing the system's own dynamics from outside.

\emph{(8) Brouwer fixed point} guarantees existence of a point the
transformation does not move---the system's own point of stillness.
Finding it requires \emph{searching} the space, a selective operation.

\emph{(9) Homology} quantifies the \emph{structure of absence} (how
many holes of each dimension). Detecting what is \emph{missing}
requires meta-knowledge about the system's own topology.

The common pattern is that each generative algebra requires the
system to operate on its own structure, audit its own losses,
classify its own modes, search its own space, or detect its own
absences---operations characterised by $\recu{k}$ (recursive
selection). A system with $k = 0$ (stagnation) cannot perform any
of them, because the content axes $(\rey{r}, \imag{i}, \exac{j})$
encode \emph{what is known}, not the capacity to interrogate what
is known.
\end{proof}

\begin{theorem}[Categorical duality as $\recu{k}$]
\label{thm:catdual}
The opposition functor $(-)^{\mathrm{op}}: \mathbf{Cat} \to
\mathbf{Cat}^{\mathrm{op}}$ is the structural realisation of
$\recu{k}$ in the category of categories. Specifically:
\begin{enumerate}[nosep]
  \item It is a meta-operation on structure (operates on morphisms, not objects)---matching $\recu{k}$'s role as the axis that acts on the other axes.
  \item It generates without cancelling (no neutral element is produced)---matching the generative character.
  \item $(-)^{\mathrm{op}} \circ (-)^{\mathrm{op}} = \mathrm{Id}$---matching the involutive character of Axiom~\ref{ax:inv}. See Remark~\ref{rem:sign} below.
\end{enumerate}
\end{theorem}

\begin{proof}
Properties (1) and (2) follow from the definition of the opposite
category: $(-)^{\mathrm{op}}$ reverses all morphisms while preserving
objects, and produces genuinely new theorems (the dual of any theorem
is a distinct theorem about the opposite structure). Property~(3) is
standard: $(\mathbf{C}^{\mathrm{op}})^{\mathrm{op}} \cong \mathbf{C}$
for any category $\mathbf{C}$.

The identification with $\recu{k}$ follows: the only axis that (a)
acts on the \emph{structure} of the other axes rather than their
content, (b) generates rather than cancels, and (c) is self-inverting,
is $\recu{k}$ by Theorem~\ref{thm:split} and Axiom~\ref{ax:inv}.
\end{proof}

\begin{remark}[Sign of self-inversion]
\label{rem:sign}
Axiom~\ref{ax:inv} states $\mathcal{K} \circ \mathcal{K} = -\mathrm{Id}$
(matching $k^2 = -1$ in quaternion algebra), while the categorical
realisation satisfies $(-)^{\mathrm{op}} \circ (-)^{\mathrm{op}} =
\mathrm{Id}$ (an involution). The essential property shared by both is
\emph{self-inversion}: applying selection twice collapses to a
determined state. The sign distinguishes two substrates: in
$\mathbb{H}$, self-inversion produces a $180^\circ$ rotation
($-\mathrm{Id}$); in $\mathbf{Cat}$, it produces the identity (exact
return). Axiom K4 captures the quaternionic version because it preserves
the full algebraic structure ($ij = k$, $jk = i$, $ki = j$), which
requires $k^2 = -1$. The categorical realisation satisfies the weaker
involutive form---it captures the generative and self-inverting character
of $\recu{k}$ but not the rotational sign.
\end{remark}

\begin{theorem}[Kernel/image as $\recu{k}$]
\label{thm:kerimg}
The rank-nullity decomposition $V = \ker(f) \oplus \ker(f)^\perp$ of
a linear map $f: V \to W$ (on finite-dimensional inner product spaces)
is the structural realisation of $\recu{k}$ in linear algebra:
\begin{enumerate}[nosep]
  \item It is a meta-operation: testing $v \in \ker(f)$ requires
    evaluating $f(v) = 0$, an examination of the map's own behaviour.
  \item It generates without cancelling: the dimensions $\dim\ker f$
    and $\operatorname{rank} f$ are global invariants invisible to any
    single evaluation $f(v)$.
  \item The complementary-projection map $P \mapsto \mathrm{Id} - P$
    is an involution: $\mathrm{Id} - (\mathrm{Id} - P) = P$.
\end{enumerate}
\end{theorem}

\begin{proof}
(1) Membership $v \in \ker(f)$ requires comparing $f(v)$ against a
criterion ($= 0$): a diagnostic on $f$'s output, not a computation
within $V$. (2) The rank-nullity identity
$\dim\ker f + \operatorname{rank} f = \dim V$ cannot be determined
from fewer than $\dim V$ linearly independent evaluations---assembling
such a spanning set is itself a selective ($\recu{k}$) operation.
(3) Immediate: $C(P) = \mathrm{Id} - P$ satisfies $C \circ C =
\mathrm{Id}$.

A system with $k = 0$ can compute $f(v)$ for given $v$ but cannot
test whether $f(v) = 0$ (no meta-diagnostic) and therefore cannot
partition $V$ into kernel and image.
\end{proof}

\begin{theorem}[Eigenvalue spectrum as $\recu{k}$]
\label{thm:spectrum}
For a normal operator $A$ on a finite-dimensional inner product space,
the spectral decomposition $A = \sum_i \lambda_i P_i$ is the structural
realisation of $\recu{k}$ in operator algebra:
\begin{enumerate}[nosep]
  \item It is a meta-operation: the characteristic polynomial
    $\det(A - \lambda I) = 0$ depends on all entries of $A$
    simultaneously.
  \item It generates without cancelling: the eigenvalues classify
    $A$'s modes (expansion, contraction, oscillation, annihilation).
  \item The adjoint map $A \mapsto A^*$ is an involution on operators
    ($A^{**} = A$) inducing conjugation $\lambda \mapsto \bar\lambda$
    on spectra.
\end{enumerate}
\end{theorem}

\begin{proof}
(1) The characteristic polynomial $p_A(\lambda) = \det(A - \lambda I)$
is a global function of $A$'s matrix entries; no single evaluation $Av$
determines any eigenvalue. (2) The spectrum $\sigma(A)$ classifies
qualitative dynamics on each eigenspace: a system limited to content
axes sees $Av$ as a vector, not as $\lambda v$ for a mode~$\lambda$.
(3) For normal $A = U\Lambda U^*$, $A^* = U\bar\Lambda U^*$, giving
$\sigma(A^*) = \overline{\sigma(A)}$, and $(A^*)^* = A$.

A system with $k = 0$ can compute $Av$ but cannot extract the
eigenvalues---it sees individual outputs without classifying the modes
that generate them.
\end{proof}

\begin{theorem}[Brouwer fixed point as $\recu{k}$]
\label{thm:brouwer}
The fixed-point operator $\mathrm{Fix}(f) = \{x \in D : f(x) = x\}$
for continuous $f: D \to D$ on compact convex $D \subset \mathbb{R}^n$
is the structural realisation of $\recu{k}$ in continuous dynamics:
\begin{enumerate}[nosep]
  \item It is a meta-operation: $f(x) = x$ compares the system's output
    with its input---a self-referential diagnostic.
  \item It generates without cancelling: Brouwer's theorem guarantees
    $\mathrm{Fix}(f) \neq \emptyset$, a topological existence result
    not determined by finitely many evaluations of $f$.
  \item The reflection $f \mapsto 2\mathrm{Id} - f$ is a
    fixed-point-preserving involution: $\mathrm{Fix}(f) = \mathrm{Fix}
    (2\mathrm{Id} - f)$, and $2\mathrm{Id} - (2\mathrm{Id} - f) = f$.
\end{enumerate}
\end{theorem}

\begin{proof}
(1) The test $f(x) \stackrel{?}{=} x$ requires simultaneous access to
$f$ and $\mathrm{Id}$---an operation on $f$'s structure, not within
its domain. (2) Brouwer's theorem follows from the Lefschetz fixed
point theorem ($L(f) \neq 0$ for contractible $D$): a global
topological invariant, not a local computation.
(3) If $f(x^*) = x^*$ then $(2x^* - f(x^*)) = x^*$, so
$x^* \in \mathrm{Fix}(2\mathrm{Id} - f)$, and conversely.
The composition $2\mathrm{Id} - (2\mathrm{Id} - f) = f$ is immediate.

A system with $k = 0$ can evaluate $f(x)$ pointwise but cannot
compare $f(x)$ with $x$ (no self-referential test) and therefore
cannot locate fixed points.
\end{proof}

\begin{theorem}[Homology as $\recu{k}$]
\label{thm:homology}
For a simplicial complex $K$, the Betti numbers
$\beta_n = \dim H_n(K) = \dim(\ker\partial_n /
\operatorname{im}\partial_{n+1})$ are the structural realisation of
$\recu{k}$ in algebraic topology:
\begin{enumerate}[nosep]
  \item It is a meta-operation: computing $H_n$ requires comparing
    $\ker\partial_n$ against $\operatorname{im}\partial_{n+1}$---an
    analysis of the boundary operators' own structure.
  \item It generates without cancelling: the Betti numbers quantify
    the ``structure of absence'' (independent $n$-cycles that do not
    bound).
  \item Poincar\'e duality for closed oriented $n$-manifolds,
    $H_k(M) \cong H_{n-k}(M)$, is a dimension-swapping involution:
    $k \mapsto n-k \mapsto k$.
\end{enumerate}
\end{theorem}

\begin{proof}
(1) The quotient $\ker\partial_n / \operatorname{im}\partial_{n+1}$
asks ``which cycles are boundaries?''---a meta-question about the
relationship between consecutive chain groups, not a computation
within any single group. (2) Betti numbers detect what is
\emph{absent}: $\beta_1$ counts independent loops, $\beta_2$ counts
independent voids. This is information about holes in $K$, not about
simplices present in $K$. (3) For closed oriented $n$-manifolds,
$\beta_k = \beta_{n-k}$, and the dimension map $k \mapsto n-k$
satisfies $n - (n-k) = k$.

A system with $k = 0$ can compute boundaries $\partial\sigma$ of
individual simplices but cannot determine which cycles bound---it sees
local boundary relations without the global analysis distinguishing
homologous from trivial cycles.
\end{proof}

\begin{remark}[Involution forms]
\label{rem:involution-forms}
Each generative algebra realises $\recu{k}$'s self-inverting character
through a different involution:
complementary projection ($P \mapsto \mathrm{Id} - P$) for
kernel/image,
operator adjunction ($A \mapsto A^*$) for eigenvalues,
dynamical reflection ($f \mapsto 2\mathrm{Id} - f$) for fixed points,
and Poincar\'e duality ($H_k \leftrightarrow H_{n-k}$) for homology.
Like Remark~\ref{rem:sign}, these are substrate-specific realisations
of the same structural property: double application returns to the
original.
\end{remark}

%======================================================================
\section{Pathologies as Axis Collapses}
\label{sec:pathologies}
%======================================================================

Each kingdom's absence produces a diagnosable pathology:

\begin{pathologybox}[title={Dogma ($\imag{i} = 0$)}]
\begin{equation}
\text{Dogma} = (\rey{r} > 0,\; \imag{i} = 0,\;
\exac{j} > 0,\; \recu{k} \geq 0)
\end{equation}
Crystallised, verified, possibly even selective---but without
potentiality. No ``what if?'' is permitted. Truth frozen.
Ice that has forgotten it was once water.
\end{pathologybox}

\begin{pathologybox}[title={Delirium ($\rey{r} = 0$)}]
\begin{equation}
\text{Delirium} = (\rey{r} = 0,\; \imag{i} > 0,\;
\exac{j} \geq 0,\; \recu{k} \geq 0)
\end{equation}
Pure potentiality without crystallisation.
Heat without container. The artist with a thousand visions who
finishes nothing.
\end{pathologybox}

\begin{pathologybox}[title={Pseudoscience ($\exac{j} = 0$)}]
\begin{equation}
\text{Pseudoscience} = (\rey{r} > 0,\; \imag{i} > 0,\;
\exac{j} = 0,\; \recu{k} \geq 0)
\end{equation}
Structure and imagination without verification.
``It sounds right'' and ``everything fits'' but nobody has
measured anything. The paper that diagnostics detects:
elegant framework, zero empirical contact.
\end{pathologybox}

\begin{pathologybox}[title={Stagnation ($\recu{k} = 0$)}]
\begin{equation}
\text{Stagnation} = (\rey{r} > 0,\; \imag{i} > 0,\;
\exac{j} > 0,\; \recu{k} = 0)
\end{equation}
Crystallised, imaginative, verified---but never pruning,
never rotating, never advancing. The committee that discusses
endlessly without deciding. Evolution without natural selection:
variation without fitness.
\end{pathologybox}

\bigskip

\begin{definition}[Liquid Zone]
\label{def:liquid}
The \textbf{liquid zone} is the region of the semantic hypercube where
all four axes are simultaneously non-zero:
\begin{equation}
\LZ = \{s \in [0,1]^4 : \rey{r} > 0 \;\land\;
\imag{i} > 0 \;\land\; \exac{j} > 0 \;\land\; \recu{k} > 0\}
\end{equation}
Phase transitions are trajectories that enter or exit $\LZ$
by collapsing or activating an axis.
\end{definition}

\subsection{Topological Properties}

\begin{proposition}[Openness]
\label{prop:open}
$\LZ$ is open in the relative topology of $H = [0,1]^4$.
\end{proposition}

\begin{proof}
For each coordinate $x \in \{r, i, j, k\}$, the set $\{s \in H : x >
0\}$ is open in $H$ (it is $H \cap \{x > 0\}$, and $\{x > 0\}$ is
open in $\mathbb{R}^4$). Since $\LZ$ is the intersection of four such
open sets, it is open.
\end{proof}

\begin{proposition}[Convexity and path-connectedness]
\label{prop:convex}
$\LZ$ is convex and hence path-connected.
\end{proposition}

\begin{proof}
Let $s_1, s_2 \in \LZ$ and $t \in [0,1]$. Then $s_t = (1-t)s_1 + t
s_2$ satisfies: each coordinate of $s_t$ is a convex combination of
two strictly positive numbers in $[0,1]$, hence strictly positive and
in $[0,1]$. Therefore $s_t \in \LZ$, and $\LZ$ is convex. Convexity
implies path-connectedness.
\end{proof}

\begin{proposition}[Disjointness of pathologies]
\label{prop:disjoint}
The four pathological regions are pairwise disjoint.
\end{proposition}

\begin{proof}
Each pathology sets a \emph{different} axis to zero while requiring at
least two others to be strictly positive. For any pair of pathologies,
at least one requires a coordinate to be positive that the other
requires to be zero (or vice versa), so their intersection is empty.
Explicitly:
$\mathcal{D}_{\text{dogma}} \cap \mathcal{D}_{\text{delirium}} =
\emptyset$ because dogma requires $r > 0$ while delirium requires
$r = 0$.
$\mathcal{D}_{\text{dogma}} \cap \mathcal{D}_{\text{pseudo}} =
\emptyset$ because dogma requires $j > 0$ while pseudoscience requires
$j = 0$.
$\mathcal{D}_{\text{dogma}} \cap \mathcal{D}_{\text{stag}} =
\emptyset$ because dogma requires $i = 0$ while stagnation requires
$i > 0$.
$\mathcal{D}_{\text{delirium}} \cap \mathcal{D}_{\text{pseudo}} =
\emptyset$ because delirium requires $r = 0$ while pseudoscience
requires $r > 0$.
$\mathcal{D}_{\text{delirium}} \cap \mathcal{D}_{\text{stag}} =
\emptyset$ because delirium requires $r = 0$ while stagnation requires
$r > 0$.
$\mathcal{D}_{\text{pseudo}} \cap \mathcal{D}_{\text{stag}} =
\emptyset$ because pseudoscience requires $j = 0$ while stagnation
requires $j > 0$.
\end{proof}

\begin{proposition}[Complementarity]
\label{prop:complement}
$\LZ \cap \mathcal{D}_{\alpha} = \emptyset$ for each pathology
$\alpha$. The liquid zone is disjoint from all pathological regions.
\end{proposition}

\begin{proof}
Every state in $\LZ$ has all four coordinates strictly positive. Every
pathological state has at least one coordinate equal to zero.
\end{proof}

%======================================================================
\section{The Quaternionic Product on \texorpdfstring{$S^3$}{S3}}
\label{sec:quaternion}
%======================================================================

The four axes $\{1, \imag{i}, \exac{j}, \recu{k}\}$ mirror the basis of
Hamilton's quaternions $\mathbb{H}$, with the multiplication rules:
\begin{equation}
\imag{i}^2 = \exac{j}^2 = \recu{k}^2 = \imag{i}\exac{j}\recu{k} = -1
\end{equation}
\begin{equation}
\imag{i}\exac{j} = \recu{k}, \quad
\exac{j}\recu{k} = \imag{i}, \quad
\recu{k}\imag{i} = \exac{j}
\end{equation}

Each rule receives a semantic interpretation:
\begin{enumerate}[nosep]
\item $\imag{i}^2 = -1$: Potentiality interacting with itself produces
  inversion. A dream that dreams itself collapses or inverts.
\item $\exac{j}^2 = -1$: Verification verifying itself without content
  produces negation. A measurement of a measurement, absent anything to
  measure, is vacuous.
\item $\recu{k}^2 = -1$: Selection selecting itself produces inversion.
  Recursion without base case is infinite regress. (Axiom~\ref{ax:inv}.)
\item $\imag{i}\exac{j} = \recu{k}$: Potentiality plus directionality
  produces selection. Imagining \emph{where to go} generates the
  capacity to choose. (Axiom~\ref{ax:gen}.)
\item $\exac{j}\recu{k} = \imag{i}$: Directionality plus selection
  produces potentiality. Constraint enables creativity.
\item $\recu{k}\imag{i} = \exac{j}$: Selection plus potentiality
  produces directionality. Choosing among possibilities creates order.
\end{enumerate}

\begin{remark}
The non-commutativity of quaternion multiplication
($\imag{i}\exac{j} \neq \exac{j}\imag{i}$) has a direct semantic
reading: imagining first and then measuring produces a different outcome
than measuring first and then imagining. The order of ontological
operations matters.
\end{remark}

\begin{theorem}[Semantic quaternionic product]
\label{thm:quat_product}
For semantic states $s_1,s_2\in\mathbb{R}^4\setminus\{0\}$,
define $\hat{s}=s/\|s\|$ and
$s_1\otimes s_2=\hat{s}_1\cdot_{\mathbb{H}}\hat{s}_2$,
where $\cdot_{\mathbb{H}}$ is Hamilton multiplication.
This product:
\begin{enumerate}[nosep]
  \item is closed on $S^3$,
  \item is associative,
  \item admits inverses,
  \item satisfies the basis relations: $ij=k$, $jk=i$, $ki=j$, $i^2=j^2=k^2=ijk=-1$,
  \item is non-commutative ($ij \neq ji$), and
  \item respects layer recovery: products of real-subspace elements remain in the real subspace.
\end{enumerate}
\end{theorem}

\begin{proof}
Write $\hat{s}_j=s_j/\|s_j\|$ for the unit normalisation.

\emph{(1) Closure.}
$\|\hat{s}_1\cdot_{\mathbb{H}}\hat{s}_2\|
=\|\hat{s}_1\|\,\|\hat{s}_2\|=1\cdot1=1$,
by the multiplicativity of the quaternion norm.

\emph{(2) Associativity.}
$\hat{s}_1\otimes(\hat{s}_2\otimes\hat{s}_3)
=\hat{s}_1\cdot_{\mathbb{H}}(\hat{s}_2\cdot_{\mathbb{H}}\hat{s}_3)
=(\hat{s}_1\cdot_{\mathbb{H}}\hat{s}_2)\cdot_{\mathbb{H}}\hat{s}_3
=(\hat{s}_1\otimes\hat{s}_2)\otimes\hat{s}_3$,
since Hamilton multiplication is associative and the normalisation
of a unit quaternion is the identity.

\emph{(3) Inverses.}
For any unit quaternion $\hat{s}$,
$\hat{s}^{-1}=\overline{\hat{s}}$ (the quaternion conjugate),
and $\hat{s}\,\overline{\hat{s}}=\|\hat{s}\|^2=1$.

\emph{(4) Basis relations.}
These are the defining relations of the quaternion algebra
$\mathbb{H}=\langle 1,i,j,k\mid i^2=j^2=k^2=ijk=-1\rangle$.

\emph{(5) Non-commutativity.}
$ij=k$ while $ji=-k$, so $ij\neq ji$.

\emph{(6) Real subspace.}
If $\Imag(\hat{s}_1)=\Imag(\hat{s}_2)=0$, then
$\hat{s}_1,\hat{s}_2\in\{\pm1\}\subset\mathbb{R}\subset\mathbb{H}$,
and their product is $\pm1$, which is real.

All six properties are additionally verified computationally at
machine precision ($\varepsilon<10^{-15}$) over $10{,}000$ random
pairs and all $72$ semantic primitives from~\citet{p2}.
\end{proof}

%======================================================================
\section{Relation to Information Theory}
\label{sec:shannon}
%======================================================================

Shannon~\citeyearpar{shannon1948} defined information along a single
axis: certainty/uncertainty. A classical bit is $\{0,1\}$ and entropy
$H = -\sum p_i \log p_i$ measures ignorance. But this captures only
\emph{resolved} information---the projection onto $\rey{r}$.

\begin{proposition}[Shannon as special case]
\label{prop:shannon}
Shannon information theory is the restriction of $\Op$ to:
\begin{equation}
\text{Shannon} = \Op\big|_{\imag{i}=0,\;\exac{j}=1,\;\recu{k}=0}
= \text{projection onto } \rey{r}
\end{equation}
The conditions are: no potentiality (all information is resolved),
full verifiability (all information is measurable), and no recursive
selection (the channel is passive).
\end{proposition}

\begin{proof}
Under the Shannon restriction: $\imag{i} = 0$ (all information is
resolved), $\exac{j} = 1$ (all information is fully verifiable),
$\recu{k} = 0$ (the channel is passive---it transmits but does not
select). Shannon's entropy operates on probability distributions over
resolved, measurable states with no selective agency. The condition
$k = 0$ explains why Shannon's theory has no concept of a channel that
\emph{chooses} what to transmit: the selective axis is collapsed.
\end{proof}

The natural extension follows the algebraic progression of number
systems:
\begin{center}
\renewcommand{\arraystretch}{1.3}
\begin{tabular}{@{}lll@{}}
\toprule
\textbf{Number system} & \textbf{Operator} &
\textbf{Information captured} \\
\midrule
$\mathbb{R}$ (reals)        & $\{0, 1\}$
  & Classical bit: presence/absence \\
$\mathbb{C}$ (complex)      & $(\rey{r}, \imag{i})$
  & Information with potentiality \\
$\mathbb{H}$ (quaternions)  & $(\rey{r}, \imag{i}, \exac{j}, \recu{k})$
  & Information with potentiality, verifiability, and selection \\
\bottomrule
\end{tabular}
\end{center}

%======================================================================
\section{Epistemic Dimensions Are Not Decoration}
\label{sec:decoration}
%======================================================================

A potential objection is that since $\Real(a)$ is always
context-invariant, the imaginary dimensions
$(\imag{i},\exac{j},\recu{k})$ are ``decoration'' that do not
affect logical reasoning. We present three arguments against this
objection.

\subsection{Operational distinguishability}

\begin{definition}[Decision function]
\label{def:decision}
A \emph{decision function} is any map $D:V\to\mathcal{A}$ from
truth values to an action space $\mathcal{A}$ that depends on the
full quaternionic truth value, not only on $\Real(a)$.
\end{definition}

\begin{proposition}
\label{prop:distinguish}
If $D$ depends on $\Imag(a)$, then there exist truth values $a,b$
with $\Real(a)=\Real(b)$ such that $D(a)\neq D(b)$.
\end{proposition}

\begin{proof}
Take $a=(r,\alpha,0,0)$ and $b=(r,0,\alpha,0)$ for any
$r\in(0,1)$ and $\alpha\neq0$. Then $\Real(a)=\Real(b)=r$
but $\Imag(a)\neq\Imag(b)$, and by assumption $D$ distinguishes them.
\end{proof}

\begin{example}[The anchor argument]
\label{ex:football}
A goal is scored: $a=(0.9,0.6,0.7,0.5)$. Both teams agree on the
fact ($\Real=0.9$). Let $q_{\mathrm{red}}$ be a context rotation
(a different ``anchor''):
\[
  q_{\mathrm{red}}\cdot a = (0.9,\;0.6,\;{-0.7},\;{-0.5}).
\]
The truth is identical. The epistemic character differs: one team
reads positive verification ($+j$) and celebrates; the other reads
negative verification ($-j$) and laments. The decision (celebrate
or lament) depends on $(i,j,k)$, not only on $r$.
\end{example}

\subsection{The panorama}

Averaging over all contexts yields a pure-real result:
\[
  \int_{\SU} q\cdot a\;d\mu(q)
  = \bigl(\Real(a),\;0,\;0,\;0\bigr),
\]
where $\mu$ is the Haar measure on $\SU$. This is the
``view from nowhere''---Boolean logic. Agents, however, do
not live in the panorama; they are situated in specific contexts.
The imaginary dimensions encode the \emph{situated perspective}
that determines behaviour.

\subsection{Epistemic quality}

Two propositions with the same truth ($\Real=0.7$) but different
epistemic profiles carry different information:
\begin{itemize}[nosep]
\item $(0.7,\;0.8,\;0.1,\;0.2)$: high potentiality, low
  verification---a \emph{hypothesis}.
\item $(0.7,\;0.1,\;0.8,\;0.2)$: low potentiality, high
  verification---an \emph{observation}.
\end{itemize}
A rational agent treats hypotheses and observations differently,
even at equal truth values. The dimensions encode
\emph{epistemic quality}, which is decision-relevant.

%======================================================================
\section{Cause--Effect Chains}
\label{sec:cause}
%======================================================================

Since $\SU$ is non-abelian, the composition of context changes is
order-dependent:
\[
  q_2\cdot(q_1\cdot a) \neq q_1\cdot(q_2\cdot a)
  \quad\text{in general.}
\]
All orderings preserve the same $\Real(a)$, but produce different
epistemic signatures $\Imag$. The process retains \emph{memory of
the path taken}: the order of perspective changes leaves a distinct
epistemic trace.

\begin{proposition}[Path dependence]
\label{prop:path}
For $n$ distinct context changes $q_1,\ldots,q_n\in\SU$, the $n!$
orderings of application to a truth value $a$ all yield the same
$\Real(a)$ but generically yield $n!$ distinct imaginary parts.
\end{proposition}

\begin{proof}
$\Real$ is invariant under every conjugation (iterated application
of Proposition~\ref{prop:action}(1)). For the imaginary part,
distinct orderings correspond to distinct elements of $\SU$ (since
the group is non-abelian), which generically produce distinct
rotations of $\Imag(a)$.
\end{proof}

This formalises the notion that ``imagining then verifying'' and
``verifying then imagining'' leave different epistemic signatures,
even when they agree on the truth. The non-commutativity is the
algebraic encoding of irreversible cause--effect chains.

%======================================================================
\section{Embedding the Prime Algebra into the G-Lattice}
\label{sec:embedding}
%======================================================================

The prime-algebraic framework of~\citet{p4} encodes each concept $C$
as a product of primes $\Phi(C)=\prod_{i=1}^{72}p_i^{s_i}$,
$s_i\in\{0,1\}$, where the 72 primitives are distributed across six
algebraic layers. The binary exponent vector
$\mathbf{s}(C)=(s_1,\ldots,s_{72})\in\{0,1\}^{72}$ turns
divisibility into componentwise $\leq$, gcd into componentwise min,
and lcm into componentwise max~\citep{birkhoff1967}.

\begin{definition}[Layer projection]
\label{def:layer_proj}
Partition the 72 indices into layer sets
$L_1,\ldots,L_6$ ($|L_1|=3$, $|L_2|=8$, $|L_3|=13$, $|L_4|=21$,
$|L_5|=23$, $|L_6|=4$). Define axis-contributing sets by the semantic
content of Table~\ref{tab:unified}:
\begin{alignat*}{2}
  R &= L_1\cup L_2\cup L_4\cup L_6 &\quad& (|R|=36),\\
  I &= L_4\cup L_6 &\quad& (|I|=25),\\
  J &= L_3\cup L_5\cup L_6 &\quad& (|J|=40),\\
  K &= L_6 &\quad& (|K|=4).
\end{alignat*}
The \emph{layer embedding} is
$e\colon\{0,1\}^{72}\to[0,1]^4$,
\[
  e(\mathbf{s}) \;=\;
  \Bigl(\;
    \frac{\sum_{i\in R}s_i}{|R|},\;
    \frac{\sum_{i\in I}s_i}{|I|},\;
    \frac{\sum_{i\in J}s_i}{|J|},\;
    \frac{\sum_{i\in K}s_i}{|K|}
  \;\Bigr).
\]
For the G-lattice domain $V=[0,1]\times[-1,1]^3$, rescale the
imaginary components: $\hat{e}=(e_r,\,2e_i-1,\,2e_j-1,\,2e_k-1)$.
\end{definition}

\begin{theorem}[Synthesis embedding]
\label{thm:embedding}
Let $A,B$ be concepts with exponent vectors $\mathbf{s}(A),\mathbf{s}(B)$.
\begin{enumerate}[nosep]
\item\label{it:mono} \textbf{Monotonicity.}
  $\Phi(A)\mid\Phi(B)$ implies
  $e(\mathbf{s}(A))\leq e(\mathbf{s}(B))$ componentwise.
\item\label{it:lcm} \textbf{Lattice compatibility.}
  $e(\mathbf{s}(\operatorname{lcm}))\geq
   \max\!\bigl(e(\mathbf{s}(A)),\,e(\mathbf{s}(B))\bigr)$
  componentwise, where $\operatorname{lcm}$
  denotes $\operatorname{lcm}(\Phi(A),\Phi(B))$.
\item\label{it:strict} \textbf{Strict synthesis.}
  Define the \emph{interiors}
  $\mathrm{Int}(A\!\mid\! B)=\{i:s_i(A)=1,\,s_i(B)=0\}$ and
  $\mathrm{Int}(B\!\mid\! A)=\{i:s_i(B)=1,\,s_i(A)=0\}$.
  The inequality in~(\ref{it:lcm}) is strict on axis $X$ if and
  only if both interiors have $X$-support:
  \[
    e(\mathbf{s}(\operatorname{lcm}))_X
    > \max(e(\mathbf{s}(A))_X,\,e(\mathbf{s}(B))_X)
    \;\;\Longleftrightarrow\;\;
    \mathrm{Int}(A\!\mid\! B)\cap X\neq\emptyset
    \;\text{and}\;
    \mathrm{Int}(B\!\mid\! A)\cap X\neq\emptyset.
  \]
\end{enumerate}
\end{theorem}

\begin{proof}
\textbf{(\ref{it:mono})}
$\Phi(A)\mid\Phi(B)$ means $s_i(A)\leq s_i(B)$ for all $i$, so
$\sum_{i\in X}s_i(A)\leq\sum_{i\in X}s_i(B)$ for every axis set $X$.

\textbf{(\ref{it:lcm})}
$s_i(\text{lcm})=\max(s_i(A),s_i(B))\geq s_i(A)$ for all~$i$, so
$\sum_{i\in X}\max(s_i(A),s_i(B))\geq\sum_{i\in X}s_i(A)$, and
similarly for~$B$. Dividing by $|X|$ gives
$e(\mathbf{s}(\text{lcm}))_X\geq\max(e_X(A),e_X(B))$.

\textbf{(\ref{it:strict})}
Let $A_X=\{i\in X:s_i(A)=1\}$ and $B_X=\{i\in X:s_i(B)=1\}$.
Then $e(\mathbf{s}(\text{lcm}))_X=|A_X\cup B_X|/|X|$ and
$\max(e_X(A),e_X(B))=\max(|A_X|,|B_X|)/|X|$.
The union strictly exceeds the larger set iff neither is a
subset of the other: $A_X\not\subseteq B_X$ and
$B_X\not\subseteq A_X$. The first holds iff
$\mathrm{Int}(A\!\mid\! B)\cap X\neq\emptyset$;
the second iff $\mathrm{Int}(B\!\mid\! A)\cap X\neq\emptyset$.
\end{proof}

\begin{remark}[Degenerate pairs]
When $\Phi(A)\mid\Phi(B)$, the interior $\mathrm{Int}(A\!\mid\! B)$
is empty and synthesis reduces to $\operatorname{lcm}=\Phi(B)$: no
axis shows strict increase because $A$ is already contained in~$B$.
Three of the fourteen duality pairs from~\citet{p4} exhibit this
degeneracy (their dependency chains are in a subset relation); the
remaining eleven all show strict increase on every axis with bilateral
interior support, as predicted by Theorem~\ref{thm:embedding}(\ref{it:strict}).
\end{remark}

%======================================================================
\section{Computational Verification}
\label{sec:computation}
%======================================================================

All results were verified computationally using NumPy in Python.
Random test cases were drawn uniformly from the appropriate domains.

\begin{table}[H]
\centering
\caption{Summary of computational verification.}
\label{tab:computation}
\begin{tabularx}{\textwidth}{@{}lrrl@{}}
\toprule
\textbf{Test} & \textbf{N} & \textbf{Violations} & \textbf{Max error} \\
\midrule
Specialisations (4 exact + 2 interpretive) & --- & 0 & exact \\
De Morgan laws (extended domain) & 50\,000 & 0 & $0$ \\
$\Real$ invariance under conjugation & 50\,000 & 0 & $5.55\times10^{-16}$ \\
$\|\Imag\|$ invariance under conjugation & 10\,000 & 0 & $6.66\times10^{-16}$ \\
Algebra automorphism & 10\,000 & 0 & $< 10^{-12}$ \\
Commutation (Boolean AND) & 10\,000 & 0 & $3.33\times10^{-16}$ \\
Non-commutation (full AND) & 10\,000 & 9\,634 & $1.16$ \\
Non-commutation (full OR) & 10\,000 & 9\,637 & $1.23$ \\
Non-commutation (NOT) & 10\,000 & 10\,000 & $2.73$ \\
Modus ponens (all contexts) & 50\,000 & 0 & exact \\
Context shift preserves truth & 50\,000 & 0 & exact \\
$\Real$ of compound formula invariant & 50\,000 & 0 & $9.99\times10^{-16}$ \\
Domain closure ($r$ in $[0,1]$) & 50\,000 & 0 & exact \\
Decision distinguishability & 50\,000 & --- & 49.6\% changed \\
Panorama ($\Imag\to 0$) & 10\,000 & --- & $\|\Imag\|<0.005$ \\
$\Real$-independence (all connectives) & 100\,000 & 0 & $0$ \\
$\Real$-independence (7 compounds) & 100\,000 & 0 & $0$ \\
$G^\sim$ tautologies (6 formulas) & 100\,000 & 0 & confirmed \\
Non-tautologies (4 formulas) & 100\,000 & 0 & confirmed \\
\midrule
\multicolumn{4}{@{}l}{\textit{Embedding (Section~\ref{sec:embedding})}} \\
Layer consistency (P1) & 4 & 0 & exact \\
Divisibility $\Rightarrow$ order (P3) & 10\,000 & 0 & exact \\
lcm $\geq$ max (P4) & 10\,000 & 0 & exact \\
Bilateral synthesis (14 pairs) & 14 & 0 & exact \\
\bottomrule
\end{tabularx}
\end{table}

Source code is available at
\texttt{scripts/\{connectives, context\_bridge, verify\_resolutions,
verify\_completeness\}.py} in the companion repository.

%======================================================================
\section{Related Work}
\label{sec:related}
%======================================================================

\begin{table}[H]
\centering
\caption{Comparison with existing algebraic structures for logic.}
\label{tab:related}
\begin{tabularx}{\textwidth}{@{}lp{3.5cm}Y@{}}
\toprule
\textbf{Structure} & \textbf{Source} & \textbf{Relation} \\
\midrule
Residuated lattice & \citet{galatos2007} & Layer~1 of our structure \\
BL-algebra & \citet{hajek1998} & Our lattice under fuzzy restriction \\
MV-algebra & \citet{cignoli2000} & Related to {\L}ukasiewicz specialisation \\
Quaternionic fuzzy sets & \citet{dai2023} & Quaternion-valued membership; no group action, no logic \\
Quantale & \citet{mulvey1986} & Complete lattice + monoid; ours uses non-abelian group \\
Quantum logic & \citet{birkhoff1936} & Projections on Hilbert spaces, not t-norms \\
Substructural logic & \citet{restall2000} & Non-commutative connectives, no group action \\
\bottomrule
\end{tabularx}
\end{table}

The closest antecedent is \citet{dai2023}, who defines
quaternionic fuzzy \emph{sets} with quaternion-valued membership
functions bounded by unit modulus. The present work differs in
three respects: (i)~we define a propositional \emph{logic} with
inference rules, soundness, and completeness, not only set
operations; (ii)~we add a group action that separates truth from
epistemic character via the commutation theorem; (iii)~our domain
is component-wise bounded ($[0,1]\times[-1,1]^3$) rather than
modulus-bounded.

More broadly, the distinguishing feature of the present work is
that the group action is not applied to the connectives (which
remain commutative) but to the \emph{truth values themselves},
with a clean invariance theorem separating objective truth from
perspectival epistemic character.

%======================================================================
\section{Conclusion}
\label{sec:conclusion}
%======================================================================

We have constructed a quaternionic operator
$\Op=\{0,1,+,\imag{i},\exac{j},\recu{k}\}$ from which six algebraic
layers are recoverable as dimensional projections---four exactly
(Boolean, G\"odel fuzzy, G\"odel--Kripke modal, and Kleene
trivalent) and two interpretively (ordinal and probabilistic)---and
shown that it supports a G-lattice---a residuated lattice equipped
with an $\SU$ group action---that genuinely unifies Boolean and fuzzy
frameworks.

The unification is not mere containment: it is mediated by a
commutation theorem (Theorem~\ref{thm:commutation}) that characterises
Boolean logic as the context-free special case. The fourth axis
$\recu{k}$ is proven irreducible (Theorem~\ref{thm:irred}), and its
presence divides the nine known algebras of opposition into
cancellative and generative classes
(Theorem~\ref{thm:split}), with formal $\recu{k}$-embeddings for
all five generative algebras
(Theorems~\ref{thm:catdual}--\ref{thm:homology}). Each axis's absence produces a
diagnosable pathology, and the region where all four axes are
simultaneously active---the liquid zone---is open, convex, and
path-connected.

Completeness reduces to $G^\sim$ (G\"odel logic with involutive
negation) via the real-part independence lemma. The imaginary
dimensions do not alter tautologies---they enrich the surrounding
structure with context sensitivity, path-dependent epistemic chains,
and operationally distinguishable epistemic dimensions. A layer-based
embedding (Theorem~\ref{thm:embedding}) connects the prime-algebraic
framework of~\citet{p4} to the G-lattice coordinates, with a bilateral
synthesis theorem characterising exactly when duality synthesis
produces strict increase on each axis.

The framework developed here is extended in three companion papers:
\citet{p11_1} investigates whether Hamilton multiplication yields a
genuinely non-commutative conjunction, proving a trilemma that rules
out such constructions on the semantic domain;
\citet{p12} develops the synthesis operation on the prime algebra and
shows that opposites generate new truth at higher algebraic layers
through an irreducibly second-order gap operation; and
\citet{p13} formalises pre-logical states---points in $V$ where the
truth axis is undefined---and proves that information emerges as
the compound concept of the pre-logical/\allowbreak logical duality,
with substrate independence demonstrated across four domains.

%======================================================================
% References
%======================================================================
\bibliographystyle{plainnat}
\begin{thebibliography}{99}

\bibitem[Belnap(1977)]{belnap1977}
Belnap, N.D. (1977).
A useful four-valued logic.
In: Dunn, J.M., Epstein, G. (eds.) \textit{Modern Uses of
Multiple-Valued Logic}, pp.~5--37. Reidel.

\bibitem[Birkhoff(1967)]{birkhoff1967}
Birkhoff, G. (1967).
\textit{Lattice Theory}, 3rd ed.
American Mathematical Society.

\bibitem[Birkhoff \& von Neumann(1936)]{birkhoff1936}
Birkhoff, G. \& von Neumann, J. (1936).
The logic of quantum mechanics.
\textit{Annals of Mathematics}, 37(4), 823--843.

\bibitem[Bou et~al.(2011)]{bou2011}
Bou, F., Esteva, F., Godo, L. \& Rodriguez, R. (2011).
On the minimum many-valued modal logic over a finite residuated
lattice.
\textit{Journal of Logic and Computation}, 21(5), 739--790.

\bibitem[Caicedo \& Rodriguez(2010)]{caicedo2010}
Caicedo, X. \& Rodriguez, R. (2010).
Standard G\"odel modal logics.
\textit{Studia Logica}, 94(2), 189--214.

\bibitem[Cignoli et~al.(2000)]{cignoli2000}
Cignoli, R., D'Ottaviano, I. \& Mundici, D. (2000).
\textit{Algebraic Foundations of Many-Valued Reasoning}.
Kluwer.

\bibitem[Dai(2023)]{dai2023}
Dai, S. (2023).
Quaternionic fuzzy sets.
\textit{Axioms}, 12(5), 490.

\bibitem[Dummett(1959)]{dummett1959}
Dummett, M. (1959).
A propositional calculus with denumerable matrix.
\textit{Journal of Symbolic Logic}, 24(2), 97--106.

\bibitem[Fitting(1991)]{fitting1991}
Fitting, M. (1991).
Many-valued modal logics.
\textit{Fundamenta Informaticae}, 15(3--4), 235--254.

\bibitem[Esteva et~al.(2000)]{esteva2000}
Esteva, F., Godo, L., H\'ajek, P. \& Navara, M. (2000).
Residuated fuzzy logics with an involutive negation.
\textit{Archive for Mathematical Logic}, 39(2), 103--124.

\bibitem[Galatos et~al.(2007)]{galatos2007}
Galatos, N., Jipsen, P., Kowalski, T. \& Ono, H. (2007).
\textit{Residuated Lattices: An Algebraic Glimpse at Substructural
Logics}.
Elsevier.

\bibitem[H\'ajek(1998)]{hajek1998}
H\'ajek, P. (1998).
\textit{Metamathematics of Fuzzy Logic}.
Kluwer.

\bibitem[Hofstadter(1979)]{hofstadter1979}
Hofstadter, D.R. (1979).
\textit{G\"odel, Escher, Bach: An Eternal Golden Braid}.
Basic Books.

\bibitem[Lounesto(2001)]{lounesto2001}
Lounesto, P. (2001).
\textit{Clifford Algebras and Spinors}, 2nd ed.
Cambridge University Press.

\bibitem[Mulvey(1986)]{mulvey1986}
Mulvey, C.J. (1986).
\&.
\textit{Rendiconti del Circolo Matematico di Palermo}, Suppl. 12,
99--104.

\bibitem[Ornelas Brand(2026a)]{p1}
Ornelas Brand, J.A. (2026).
Triadic Neurosymbolic Engine: Prime Factorization as a Neurosymbolic
Bridge: Projecting Continuous Embeddings into Discrete Algebraic Space
for Deterministic Verification.
\textit{Zenodo}.
\href{https://doi.org/10.5281/zenodo.19205804}{doi:10.5281/zenodo.19205804}

\bibitem[Ornelas Brand(2026b)]{p2}
Ornelas Brand, J.A. (2026).
End-to-End Prime Factorization in a Generative Language Model:
Emergent Algebraic Semantics from Joint Training.
\textit{Zenodo}.
\href{https://doi.org/10.5281/zenodo.19206544}{doi:10.5281/zenodo.19206544}

\bibitem[Ornelas Brand(2026c)]{p4}
Ornelas Brand, J.A. (2026).
Toward a Layered Ontology of Semantic Duality: 14 Candidate Dualities
Across 6 Algebraic Layers with Preliminary Domain and Neural Evidence
(v1.0).
\textit{Zenodo}.
\href{https://doi.org/10.5281/zenodo.19375166}{doi:10.5281/zenodo.19375166}

\bibitem[Ornelas Brand(2026d)]{p11_1}
Ornelas Brand, J.A. (2026).
Toward a Non-Commutative Residuated Lattice from Quaternion
Multiplication (v0.1.0).
\textit{Zenodo}.
\href{https://doi.org/10.5281/zenodo.19561406}{doi:10.5281/zenodo.19561406}

\bibitem[Ornelas Brand(2026e)]{p12}
Ornelas Brand, J.A. (2026).
Duality Synthesis in Quaternionic Logic: How Opposites Generate Truth (v0.1.0).
\textit{Zenodo}.
\href{https://doi.org/10.5281/zenodo.19561633}{doi:10.5281/zenodo.19561633}

\bibitem[Ornelas Brand(2026f)]{p13}
Ornelas Brand, J.A. (2026).
Pre-Logical States and the Birth of Information: How Reasoning
Operates Before Truth Exists (v0.1.0).
\textit{Zenodo}.
\href{https://doi.org/10.5281/zenodo.19561721}{doi:10.5281/zenodo.19561721}

\bibitem[Porteous(1995)]{porteous1995}
Porteous, I.R. (1995).
\textit{Clifford Algebras and the Classical Groups}.
Cambridge University Press.

\bibitem[Restall(2000)]{restall2000}
Restall, G. (2000).
\textit{An Introduction to Substructural Logics}.
Routledge.

\bibitem[Shannon(1948)]{shannon1948}
Shannon, C.E. (1948).
A mathematical theory of communication.
\textit{The Bell System Technical Journal}, 27(3), 379--423.

\end{thebibliography}

\end{document}
